\documentclass[11pt,a4paper,oneside,openany]{book}
\usepackage[T1]{fontenc}
\usepackage[utf8]{inputenc}
\usepackage[english]{babel}
\usepackage{lmodern,amsmath,amssymb,amsthm,mathtools,bm}
\usepackage{booktabs,longtable,geometry,graphicx,microtype,enumitem}
\usepackage{hyperref,xcolor,fancyhdr}
\geometry{margin=2.55cm,headheight=14pt}
\hypersetup{colorlinks=true,linkcolor=blue!55!black,citecolor=blue!55!black,urlcolor=blue!55!black,pdftitle={From Ordinary Language to Computational Neuroscience},pdfauthor={CON2EPHYS research notebook}}
\numberwithin{equation}{chapter}
\newtheorem{definition}{Definition}[chapter]
\newtheorem{proposition}[definition]{Proposition}
\newtheorem{theorem}[definition]{Theorem}
\newtheorem{remark}[definition]{Remark}
\theoremstyle{definition}
\newtheorem{workedexample}[definition]{Worked example}
\newtheorem{readergoal}[definition]{Reading goal}
\newenvironment{exercises}{\par\medskip\noindent\textbf{Exercises}\begin{enumerate}[leftmargin=2em,itemsep=0.18em]}{\end{enumerate}\medskip}
\newcommand{\E}{\mathbb E}
\newcommand{\Var}{\operatorname{Var}}
\newcommand{\Cov}{\operatorname{Cov}}
\newcommand{\R}{\mathbb R}
\newcommand{\C}{\mathbb C}
\newcommand{\Prob}{\mathbb P}
\newcommand{\diag}{\operatorname{diag}}
\newcommand{\rank}{\operatorname{rank}}
\pagestyle{fancy}\fancyhf{}\fancyhead[L]{\small CON$^2$EPHYS Mathematics}\fancyhead[R]{\small\nouppercase{\leftmark}}\fancyfoot[C]{\thepage}
\setlength{\parskip}{0.25em}
\title{\Huge From Ordinary Language to\\ Computational Neuroscience\\[1em]\Large A Mathematical Foundation for Electrophysiological Analysis\\[0.5em]\normalsize A modular companion to CON$^2$EPHYS \texttt{arsenal/path}}
\author{CON$^2$EPHYS Research Manuscript}
\date{Expanded foundational edition --- October 2026\\\small Independent exposition inspired by Halmos and Mendelson}
\begin{document}
\frontmatter
\maketitle
\chapter*{Preface}\addcontentsline{toc}{chapter}{Preface}
This book begins with statements in ordinary English, not with calculus, matrix notation, or a list of methods. Its reader need not know what a set, function, derivative, matrix, probability, or oscillation is. We will name each mathematical object before asking the reader to manipulate it, and we will distinguish statements about abstract models from statements about nervous systems. The purpose is not to make mathematics superficial: it is to make the deepest methods accessible by building their prerequisites instead of assuming them.

We follow two complementary pedagogical inspirations. Halmos develops the set-theoretic language of mathematics through informal but precise exposition; Mendelson connects logic, sets, and successive number systems to the foundations of analysis. Neither book is copied here. Their themes guide an original scientific text culminating in signal processing, neuronal point processes, statistical inference, and population dynamics.

Every chapter includes formal mathematics, an electrophysiological interpretation, limitations, and exercises. A reader should practice proving elementary results as well as simulating advanced algorithms. Source-data quality, estimand definitions, and null hypotheses receive the same attention as derivations. The linked \texttt{arsenal/path} index is currently provisional until the exact GitHub tree can be verified.
\tableofcontents
\chapter{How to read this book}
The intended reader begins with ordinary English and curiosity about neural signals. Part I requires no mathematical notation beyond what it introduces. Parts II and III build number systems, limits, and finite-dimensional algebra before the book assumes calculus, probability, or complex phase. Later chapters use the same format: meaning in words, formal definition, calculation or derivation, a neuroscience connection, a failure mode, and exercises. A novice may work every exercise; an experienced reader can use the cross-reference path below to fill only missing prerequisites.

\section*{Suggested prerequisite chains}
\begin{description}[leftmargin=2cm,style=nextline]
\item[Hilbert and PAC] English $\to$ functions $\to$ complex numbers $\to$ sequences/integrals $\to$ Fourier $\to$ analytic signal $\to$ circular inference.
\item[PCA and manifolds] Relations $\to$ functions $\to$ real numbers $\to$ vector spaces $\to$ covariance $\to$ SVD/PCA $\to$ dynamical systems.
\item[Spike GLMs] Sets $\to$ functions $\to$ counting $\to$ probability $\to$ conditional probability $\to$ Poisson processes $\to$ likelihood and GLM.
\item[Granger analysis] Covariance $\to$ stochastic processes $\to$ cross spectra $\to$ multivariate linear prediction $\to$ null models and validation.
\item[Current source density] Coordinates $\to$ derivatives $\to$ multivariable calculus $\to$ forward measurement physics $\to$ regularized inverse methods.
\end{description}

\section*{Technical vocabulary}
An \emph{estimand} is the quantity the analysis intends to estimate; an \emph{estimator} is a data-dependent rule for estimating it. A \emph{model} is a set of assumed mathematical relationships; a \emph{measurement} is an observation through an instrument. \emph{Association} is not automatically \emph{causation}. \emph{Prediction} need not identify an intervention effect. The word ``significant'' should be accompanied by the hypothesis and statistical criterion being used.

\section*{About the two foundational books}
The structure draws inspiration from Halmos' conversational construction of sets and relations \cite{HalmosNaive}, and Mendelson's progression from logic and set theory to the number systems and real-number completeness \cite{MendelsonNumbers}. The explanations, examples, and proofs here are written afresh for CON$^2$EPHYS and are not reproduced sections of either text.
\mainmatter
\part{The language of mathematics}
\chapter{From ordinary English to mathematical statements}\label{ch:english}
\begin{readergoal}Read, question, and translate statements before manipulating symbols. No mathematical prerequisite is assumed.\end{readergoal}
\section{Why begin with sentences?}
A scientist looks at an experiment and makes claims: ``this animal ran faster,'' ``these spikes followed the stimulus,'' or ``these electrodes recorded the same rhythm.'' At first, the claims are expressed in ordinary language. Mathematics begins when we decide exactly what would make a claim true, how it could be false, and what follows from it. The symbols are a shorthand for decisions already made in words. Learning a formula without learning its grammar is like memorizing a sentence in a language one cannot read.

Consider ``every recorded neuron fired.'' Does this refer to every neuron in the brain, every electrode channel, or every neuron retained after spike sorting? The statement is incomplete until we specify its domain of discourse. Consider ``some neurons were silent.'' The word \emph{some} means \emph{at least one}; it does not exclude the possibility that many neurons were silent. The word \emph{or} in standard mathematics ordinarily permits both alternatives. When a manuscript says ``the signal is oscillatory,'' we must also decide what counts as an oscillation. A label is not yet an operational test.

\section{Statements, questions, definitions, and hypotheses}
A \emph{proposition} is a declarative statement to which a truth value can be assigned, given a specified interpretation. ``Two is even'' is a proposition. ``Is two even?'' is a question, while ``Let $n$ be the number of spikes'' is an instruction introducing notation. ``For every $t$ in the recording interval, $x(t)>0$'' becomes a proposition once $x$ and the recording interval are fixed. An \emph{open sentence} such as ``$n$ is even'' has a free variable and does not yet assert anything about a particular number.

A \emph{definition} establishes how a term or symbol will be used; a \emph{hypothesis} asserts a fact that may or may not hold; a \emph{theorem} is a proposition established from previously stated assumptions by valid reasoning. In applied science we must further distinguish a mathematical theorem, which follows within a model, from an empirical assertion about the world, which requires observations.

\section{Translating ordinary English}
Let $U$ be the set of units retained after quality control. Introduce the predicate $F(u)$ for ``unit $u$ fired in a designated 100-ms window.'' Then ``every retained unit fired'' becomes
\[ \forall u\in U,\ F(u). \]
``At least one retained unit did not fire'' becomes $\exists u\in U$ such that $\neg F(u)$. These are logical negations of one another. ``No unit fired'' instead means $\forall u\in U,\neg F(u)$. Notice the difference between a negative statement about \emph{all} units and a statement that \emph{not all} units fired.

In data analysis, ``higher gamma amplitude when theta phase is near its peak'' must be turned into an explicit comparison: define phase bins, summarize amplitudes, select a test statistic, and specify a null hypothesis. This transition from sentence to mathematical object is the beginning of sound methodology.

\section{A claim is only as good as its scope}
Suppose 40 of 50 detected units increase firing rate during a task. It is legitimate to say ``80 percent of these detected units increased.'' It is not automatically legitimate to say ``80 percent of neurons in this brain area increased.'' Electrode sampling, sorting errors, and selection rules change the population to which a claim refers. We will repeatedly separate \emph{a set of observed objects}, \emph{a mathematical model}, and \emph{a physiological population}.

\begin{workedexample}[The ambiguity of ``only'']
Compare ``only neurons in CA1 show a theta rhythm'' with ``neurons in CA1 show only a theta rhythm.'' The first restricts the location of units exhibiting theta; the second restricts the rhythms observed in CA1 units. Write down what observation would contradict each statement. The two claims can disagree even on the same recording.
\end{workedexample}
\begin{exercises}
\item Translate ``not every trial contains a ripple'' into a quantified sentence and simplify its negation.
\item Explain why ``all spikes occur in bursts'' can be vacuously true for an empty spike train.
\item Identify the observational unit in each claim: neuron, trial, event, animal, and session.
\item Write one empirically falsifiable claim about coherence and one purely mathematical definition of coherence.
\end{exercises}
\chapter{Logic: the grammar of proof}\label{ch:logic}
\begin{readergoal}Learn the connectives, truth conditions, quantifiers, and the difference between logical implication and scientific causation.\end{readergoal}
\section{And, or, and not}
Let $P$ mean ``the animal is running'' and $Q$ mean ``theta power exceeds a specified threshold.'' The conjunction $P\land Q$ is true precisely when both claims are true. The disjunction $P\lor Q$ is false only when both are false; in mathematical usage it is \emph{inclusive}. Negation $\neg P$ reverses the truth value of $P$. The entire truth table is small enough to check directly:
\[\begin{array}{cc|ccc}P&Q&P\land Q&P\lor Q&P\Rightarrow Q\\\hline
T&T&T&T&T\\T&F&F&T&F\\F&T&F&T&T\\F&F&F&F&T\end{array}\]
The conditional $P\Rightarrow Q$ is false only when $P$ holds but $Q$ does not. This is the \emph{material implication}; it is not an assertion of a mechanism or a causal influence. Indeed, a false antecedent makes the material implication true without producing any biological event.

\section{Converse, inverse, and contrapositive}
From $P\Rightarrow Q$ one may infer the logically equivalent contrapositive $\neg Q\Rightarrow\neg P$. The converse $Q\Rightarrow P$ and inverse $\neg P\Rightarrow\neg Q$ are different claims. A low firing rate might be observed during inactivity, but inactivity need not be the only circumstance producing that rate. We may write $P\Leftrightarrow Q$ for a \emph{biconditional}, meaning that both implications hold.

De Morgan's rules supply a useful bridge to data selection: $\neg(P\land Q)\Leftrightarrow(\neg P\lor\neg Q)$ and $\neg(P\lor Q)\Leftrightarrow(\neg P\land\neg Q)$. Thus the complement of trials satisfying both inclusion criteria is the collection failing at least one criterion.

\section{Quantifiers and their order}
The phrase ``for every'' is denoted $\forall$, and ``there exists'' by $\exists$. These operations are not interchangeable: $\forall i\,\exists t\,F(i,t)$ says each neuron fires at some time, possibly a different time for every neuron. The stronger statement $\exists t\,\forall i\,F(i,t)$ claims there is one time at which every neuron fires. Neither form alone specifies simultaneity to a millisecond: that depends on how $t$ and $F$ are defined.

Negation propagates through quantifiers by the rules
\[\neg\forall x\,P(x)\ \Leftrightarrow\ \exists x\,\neg P(x),\qquad
\neg\exists x\,P(x)\ \Leftrightarrow\ \forall x\,\neg P(x).\]
For a finite dataset, verifying a universal claim requires checking all included cases; refuting it requires only one counterexample.

\section{Necessary and sufficient conditions}
If $P\Rightarrow Q$, then $P$ is \emph{sufficient} for $Q$, while $Q$ is \emph{necessary} for $P$. An example: if a matrix $A$ is invertible, its columns are linearly independent. In the square finite-dimensional case, the converse also holds; the statements are equivalent. Contrast this with ``if the signal is a sinusoid, it has a spectral peak,'' whose converse may fail for many reasons. Whenever an algorithm claims to ``detect oscillations,'' ask whether its acceptance rule is necessary, sufficient, both, or neither for the physiological interpretation.

\begin{workedexample}[A null hypothesis in plain language]
Suppose $H_0$ states that the temporal positions of spikes inside specified jitter windows carry no additional information beyond their counts. Rejecting $H_0$ is evidence against this particular conditional model. It does \emph{not} establish that a synapse connects the neurons, or that all alternative null models would also be rejected. Compare \cite{Amarasingham2012}.
\end{workedexample}
\begin{exercises}
\item Construct a complete truth table for $(P\Rightarrow Q)\Leftrightarrow(\neg Q\Rightarrow\neg P)$.
\item Negate ``for every frequency there exists a channel with nonzero power.''
\item Give a counterexample to the inference from $P\Rightarrow Q$ and $Q$ to $P$.
\item Explain why statistical dependence is not equivalent to causal influence.
\end{exercises}
\chapter{Definitions, proof, and mathematical maturity}\label{ch:proof}
\begin{readergoal}Learn to construct and audit arguments, identify hidden hypotheses, and write counterexamples.\end{readergoal}
\section{The anatomy of an argument}
A proof begins with an exact statement and a list of assumptions. It proceeds through steps justified by a definition, an axiom, or an earlier result. The conclusion is valid \emph{conditional on those assumptions}. Mathematical evidence and empirical evidence answer different questions: a covariance matrix is positive semidefinite by construction, but a measured covariance need not describe only neuronal interactions.

\section{Direct proof and proof by contrapositive}
We say an integer $n$ is even if $n=2k$ for some integer $k$. To prove that the sum of two even integers is even, write $a=2k$ and $b=2m$. Then $a+b=2(k+m)$; since $k+m$ is an integer, the sum is even. Notice the explicit use of the definition. An assertion such as ``the result is obvious'' hides exactly the step a beginner may need.

For the contrapositive method, consider: if $n^2$ is even, then $n$ is even. Prove instead that odd $n=2k+1$ implies $n^2=4k(k+1)+1$, which is odd. This avoids assuming the conclusion.

\section{Contradiction, induction, and counterexamples}
Proof by contradiction begins with the negation of the proposed conclusion, then derives an impossibility. The classic irrationality argument for $\sqrt2$ assumes $\sqrt2=a/b$ with coprime integers $a,b$, derives that both must be even, and contradicts coprimality. We shall revisit this when constructing the real numbers.

A single valid counterexample refutes a universal proposition. To refute ``every high-power spectral peak is neuronal,'' one can simulate a sinusoidal electrical artifact. To refute ``correlation implies coupling,'' one can simulate two independently noisy signals influenced by a common source. These examples do not show that neuronal oscillations or coupling never occur; they refute the unrestricted inference.

\section{Existence, uniqueness, and identifiability}
Three questions recur throughout this book: Does the object exist? Is it uniquely determined? Can we infer it from the observations? The first two are mathematical; the third also depends on a measurement model. For example, different neuronal current distributions may produce the same measured extracellular potential, so a valid forward model does not guarantee a unique inverse solution.

\begin{proposition}If $A$ and $B$ are subsets of a fixed universe $U$, then $A=B$ if and only if both $A\subseteq B$ and $B\subseteq A$.\end{proposition}
\begin{proof}If $A=B$, either inclusion follows immediately. Conversely, the first inclusion implies every element of $A$ belongs to $B$, and the second implies every element of $B$ belongs to $A$. The two sets therefore have exactly the same elements.\end{proof}
\begin{exercises}
\item Prove that the intersection of two subsets is contained in each subset.
\item Give two distinct functions with the same value at zero.
\item Distinguish ``the algorithm converges'' from ``the result estimates the intended biological quantity.''
\item Write a three-line counterexample to ``a zero Pearson correlation implies independence.''
\end{exercises}
\chapter{Sets and the operations that create them}\label{ch:sets}
\begin{readergoal}Use membership, inclusion, unions, intersections, set differences, power sets, and the idea of an admissible collection.\end{readergoal}
\section{From a collection to a set}
A set is a mathematical collection whose members are determined by a membership rule. If $U=\{a,b,c\}$, then $a\in U$ and $d\notin U$. In the ordinary language of mathematics, $a\in U$ means that $a$ \emph{belongs to} $U$; $A\subseteq U$ means every element of $A$ belongs to $U$. These statements differ: $\{a\}\subseteq U$, but typically $\{a\}\notin U$. The empty set $\varnothing$ has no members; $\{\varnothing\}$ has one member, which happens to be the empty set.

Two sets are equal exactly when they have identical members (the extensional criterion). Mathematical set construction must take place within an explicit foundational framework; one cannot freely declare every grammatical property to specify a set of all possible objects. The unrestricted collection ``all sets that are not members of themselves'' generates Russell's paradox. Following the accessible axiomatic spirit of \cite{HalmosNaive,MendelsonNumbers}, we allow subsets defined by conditions on an already established set and use conventional constructions for other sets.

\section{Set operations from ordinary English}
Let $A$ be the trials in which the animal runs, and $B$ those with a detected theta peak. Then $A\cap B$ consists of trials satisfying \emph{both} criteria, while $A\cup B$ contains those satisfying \emph{at least one}. Set difference $A\setminus B$ contains trials that are in $A$ but not $B$. Relative complement $A^c$ is meaningful only after choosing the universe $U$: here $A^c=U\setminus A$. We must never conflate the universe of trials in one session with all possible trials.

De Morgan's identities are
\[(A\cup B)^c=A^c\cap B^c,\qquad(A\cap B)^c=A^c\cup B^c.\]
For $U=\{1,2,3,4,5\}$, $A=\{1,3,5\}$, and $B=\{3,4\}$, we have $A\cap B=\{3\}$ and $A\cup B=\{1,3,4,5\}$. Work through the identities element by element before trusting the symbols.

\section{Families, indexing, and power sets}
A \emph{family} $(A_i)_{i\in I}$ is an indexed collection of sets. The union $\bigcup_{i\in I}A_i$ collects elements present in at least one $A_i$; the intersection $\bigcap_{i\in I}A_i$ collects elements present in every $A_i$. The \emph{power set} $\mathcal P(A)$ is the set of all subsets of $A$. For $A=\{u,v\}$ it is $\{\varnothing,\{u\},\{v\},\{u,v\}\}$. If $A$ has $n$ elements, its power set contains $2^n$ elements: each original element can independently be included or excluded.

\section{The significance for scientific data}
Many preprocessing operations change the set of data included in analysis. If $U$ is the set of detected units, then $U_{\rm quality}$ might retain only units meeting refractory-violation and signal-to-noise thresholds. Estimating a correlation on $U_{\rm quality}$ is not the same question as estimating it on $U$. Record the inclusion criteria as mathematical predicates; this makes sampling assumptions auditable.

\begin{workedexample}[Counting selections]
Suppose four electrodes pass quality control. A subset represents the electrodes chosen for a multichannel analysis; there are $2^4=16$ possible subsets, including the empty selection. Requiring at least one electrode yields 15. Requiring exactly two yields $\binom42=6$. This distinction prepares us for combinatorics and multiple-comparison correction.
\end{workedexample}
\begin{exercises}
\item Prove $A\cap(B\cup C)=(A\cap B)\cup(A\cap C)$ from element membership.
\item Determine whether $\varnothing\in\{\varnothing\}$ and whether $\varnothing\subseteq\{\varnothing\}$.
\item Describe the difference between a population of neurons and its sampled subset.
\item Given three data-exclusion rules, express the accepted trials with set operations.
\end{exercises}
\chapter{Ordered pairs, products, and relations}\label{ch:relations}
\begin{readergoal}Understand how ordered observations lead to relations and why correlation, equality, and ordering are mathematically different relations.\end{readergoal}
\section{The order of information matters}
The collection $\{2,7\}$ does not specify which element came first. An \emph{ordered pair} $(2,7)$ does. Formally, an ordered-pair construction must satisfy
\[(a,b)=(c,d)\quad\Longleftrightarrow\quad a=c\ \text{and}\ b=d.\]
One set-theoretic realization is $(a,b)=\{\{a\},\{a,b\}\}$, although the particular encoding is much less important than the stated property. A spike observation might be the pair $(t,u)$ of timestamp and unit identifier; reversing it gives a different type of object.

The \emph{Cartesian product} $A\times B$ contains all ordered pairs $(a,b)$ with $a\in A$ and $b\in B$. If $A$ has $m$ elements and $B$ has $n$, then $A\times B$ has $mn$ elements. More generally, $A_1\times\cdots\times A_k$ defines ordered tuples. A multichannel recording may be indexed by channel $c$, time sample $n$, and trial $r$: its index space is $C\times T\times R$.

\section{What is a relation?}
A \emph{binary relation} from $A$ to $B$ is a subset $R\subseteq A\times B$. Writing $aRb$ is shorthand for $(a,b)\in R$. One relation on the integers is $a\le b$; another is ``$a$ divides $b$.'' In neuroscience, an empirically defined relation may connect two neurons if their estimated spike-count correlation is above a threshold. This makes a graph edge, not necessarily a synapse.

A relation $R\subseteq A\times A$ is \emph{reflexive} if $aRa$ for every $a$, \emph{symmetric} if $aRb$ implies $bRa$, and \emph{transitive} if $aRb$ and $bRc$ imply $aRc$. A relation satisfying all three is an \emph{equivalence relation}; it partitions $A$ into disjoint equivalence classes. A \emph{partial order} is reflexive, antisymmetric, and transitive. Equality is an equivalence relation; subset inclusion is a partial order; temporal precedence is an order relation that requires a convention for simultaneous events.

\section{Composing and inverting relations}
For $R\subseteq A\times B$ and $S\subseteq B\times C$, define
\[S\circ R=\{(a,c):\exists b\in B,\ (a,b)\in R\ \land\ (b,c)\in S\}.\]
The inverse relation is $R^{-1}=\{(b,a):(a,b)\in R\}$. Relation inversion is not yet the same as inversion of a function. For instance, a relation recording which neurons are active during which trials can be reversed to ask which trials contain activity from a neuron.

\begin{workedexample}[Why graph clustering needs care]
Let $uRv$ mean ``estimated Pearson correlation between $u$ and $v$ exceeds 0.5.'' The rule is symmetric whenever correlations are defined, and self-pairs have correlation one for nonconstant variables, but it need not be transitive: neuron 1 may correlate strongly with 2, 2 with 3, and 1 weakly with 3. Thus its connected components are not the same as equivalence classes generated by the direct threshold relation.
\end{workedexample}
\begin{exercises}
\item List $A\times B$ for $A=\{a,b\}$, $B=\{0,1\}$, and explain why $A\times B\ne B\times A$ generally.
\item Construct a reflexive symmetric relation that is not transitive.
\item Prove the inverse of an equivalence relation is the same relation.
\item Interpret a directed graph edge as a relation; explain why its direction does not prove causal direction.
\end{exercises}
\chapter{Functions: inputs, outputs, and transformations}\label{ch:functions}
\begin{readergoal}Distinguish relations from functions; work with domain, codomain, image, inverse, and composition.\end{readergoal}
\section{A function is a special relation}
A function $f:A\to B$ assigns to each $a\in A$ exactly one value $f(a)\in B$. It may assign the same output to many inputs. The set $A$ is its \emph{domain}, $B$ its \emph{codomain}, and $f(A)=\{f(a):a\in A\}$ its image or range. A relation that assigns two distinct outputs to the same input is not a function. These distinctions matter when saying ``the neural signal determines the behavior'': an observed relation between activity and behavior is not automatically a deterministic function.

A function is \emph{injective} when different inputs have different outputs, \emph{surjective} when every point in the codomain occurs as an output, and \emph{bijective} when both conditions hold. Only a bijection $f:A\to B$ has a genuine inverse function $f^{-1}:B\to A$. Every relation has an inverse relation, but an arbitrary function does not have an inverse function.

\section{Compositions and transformations}
Given $f:A\to B$ and $g:B\to C$, their composition is $g\circ f:A\to C$, where $(g\circ f)(a)=g(f(a))$. Composition generally does not commute. In electrophysiology, one might take raw voltage $x$, apply a filter $F$, and then a nonlinear amplitude transform $H$: $H(F(x))$. Reversing the order may be meaningless or produce a different result.

\section{Sequences are functions}
A sequence $(x_n)_{n\in\mathbb N}$ is a function from an index set of natural numbers into a set of values. A finite sampled signal is a function $x:\{0,\ldots,N-1\}\to\mathbb R$. A continuously indexed signal is a function $x:I\to\mathbb R$ for some time interval $I$. A multichannel signal is $x:I\to\mathbb R^C$. The difference between continuous and sampled time is not a notational accident: sampling introduces an observation process and potentially destroys information.

\section{Functions of functions}
A \emph{functional} maps functions to numbers. For example, energy over an interval can be defined as $E[x]=\int_I |x(t)|^2\,dt$. A \emph{linear operator} maps functions to functions while respecting addition and scalar multiplication. Filtering, differentiation, and Fourier transformation are operators under suitable assumptions. These ideas will become precise after we construct number systems, vector spaces, and limits.

\begin{workedexample}[A map that loses information]
The square map $q:\mathbb R\to\mathbb R$, $q(x)=x^2$, is not injective: $q(-2)=q(2)=4$. Its inverse relation associates $4$ with both $2$ and $-2$. Restricting the domain to $[0,\infty)$ and codomain to $[0,\infty)$ makes $q$ bijective and its inverse is the nonnegative square root. Signal power similarly discards phase and sign information: it cannot reconstruct the original signal by itself.
\end{workedexample}
\begin{exercises}
\item Decide whether $f(x)=x^3$ from $\mathbb R$ to $\mathbb R$ is bijective.
\item State domain and codomain for a histogram of spike counts in 10-ms bins.
\item Find $f\circ g$ and $g\circ f$ for $f(x)=x+1$ and $g(x)=x^2$.
\item Explain why a spectrogram generally does not contain sufficient information to uniquely recover a waveform without additional phase assumptions.
\end{exercises}
\part{The construction of number systems}
\chapter{Natural numbers, counting, and induction}\label{ch:naturals}
\begin{readergoal}See how counting can be built from basic rules, rather than accepted as an unexplained operation.\end{readergoal}
\section{Counting precedes measurement}
When we say there were seven action potentials, we count discrete events. Counting does not require time units, voltage calibration, or a model of firing probability. The natural numbers are the first formal setting for such statements. We adopt $\mathbb N_0=\{0,1,2,\ldots\}$, reserving $\mathbb N_+=\{1,2,\ldots\}$. Authors differ about including zero, so every technical treatment should make its convention explicit.

\section{Successors and the Peano viewpoint}
Imagine a set $N$ with a distinguished element $0$ and successor operation $S:N\to N$. The basic idea is that zero is not the successor of anything; distinct inputs have distinct successors; and any subset containing zero and closed under the successor operation contains all of $N$. The final clause is the induction principle. These structural properties establish a unique natural-number system up to isomorphism, under the usual axiomatic background \cite{MendelsonNumbers}.

Define addition recursively by $n+0=n$ and $n+S(m)=S(n+m)$. Define multiplication through $n\cdot0=0$ and $n\cdot S(m)=n\cdot m+n$. These are \emph{definitions}, while commutativity, associativity, and distributivity are derived properties. We ordinarily write $n+1$ rather than $S(n)$ after connecting the symbols.

\section{Induction as a way to prove infinite families of claims}
Suppose we wish to prove $\sum_{k=1}^n k=n(n+1)/2$ for all positive integers $n$. The base case $n=1$ holds. Assume it is true for $n$. Then
\[\sum_{k=1}^{n+1}k=\frac{n(n+1)}2+(n+1)=\frac{(n+1)(n+2)}2.\]
Therefore the claim holds for $n+1$; induction establishes every case. Induction is a proof principle, not an empirical inference from observing a finite number of examples.

\section{Finite sets and combinatorics}
Two finite sets have the same cardinality if there exists a bijection between them. The number of length-$n$ binary vectors is $2^n$; the number of ways to choose $k$ distinct objects from $n$ without order is $\binom nk=n!/[k!(n-k)!]$. If 20 channels are inspected in all unordered pairs, the number of pairs is $\binom{20}{2}=190$. This count matters when reporting a family of connectivity tests.

\begin{workedexample}[Discrete time windows]
A recording sampled at $f_s=1000$ samples/s for 10 seconds has $N=f_sT=10000$ samples if we choose the index set $\{0,\ldots,9999\}$. The sample at index $n$ corresponds to $t_n=n/f_s$, so the last recorded time is $9.999$ s, not 10 s. Distinguish the \emph{duration of acquisition} from the \emph{timestamp of its final sample}.
\end{workedexample}
\begin{exercises}
\item Prove by induction that $1+3+\cdots+(2n-1)=n^2$.
\item Count unordered neuron pairs from $N=80$ recorded units.
\item Show that a finite sequence of $N$ entries is a function on $\{0,\ldots,N-1\}$.
\item Explain why mathematical induction and statistical generalization from trials are different operations.
\end{exercises}
\chapter{Integers: extending counting to differences}\label{ch:integers}
\begin{readergoal}Understand negative quantities through equivalence classes and learn what a legitimate algebraic extension must preserve.\end{readergoal}
\section{Why natural numbers are not enough}
A count cannot be negative, but a \emph{difference between counts} can be. We need numbers representing expressions like $3-5$ without insisting that subtraction remain inside the natural numbers. Define a relation on pairs $(a,b)\in\mathbb N_0^2$ by
\[(a,b)\sim(c,d)\quad\Longleftrightarrow\quad a+d=b+c.\]
This is an equivalence relation; its classes represent formal differences $a-b$. For example, $(3,5)\sim(0,2)$ because $3+2=5+0$. We call the set of equivalence classes $\mathbb Z$.

\section{Operations on classes}
Write $[(a,b)]$ for the class of $(a,b)$. Addition can be defined by
\[[(a,b)]+[(c,d)]=[(a+c,b+d)],\]
and multiplication by
\[[(a,b)]\cdot[(c,d)]=[(ac+bd,ad+bc)].\]
A good definition must be \emph{well defined}: choosing different representatives of the same classes must not change the resulting class. This issue appears again in quotient spaces and latent-variable models; an output must not depend on an arbitrary representation of the same underlying object.

\section{Order and absolute value}
Integers inherit addition, multiplication, and a total order. Negative numbers permit signed offsets and differences, but not all integers possess multiplicative inverses within $\mathbb Z$. The absolute value $|n|$ removes sign: $|n|=n$ when $n\ge0$, and $|n|=-n$ otherwise. It satisfies the triangle inequality $|a+b|\le|a|+|b|$.

\section{A neuroscience example}
Spike-count differences $D=N_{\rm post}-N_{\rm pre}$ may be negative although each count is nonnegative. However, raw subtraction is sensitive to the lengths of the counting intervals. If the intervals differ, compare normalized rates $N/T$ rather than interpreting the count difference uncritically. Mathematical legitimacy does not guarantee scientific comparability.

\begin{exercises}
\item Verify symmetry and transitivity of the equivalence relation on ordered pairs.
\item Show that $[(a,b)]+[(b,a)]$ is the integer zero.
\item Explain why subtraction of two nonnegative power estimates need not be nonnegative.
\end{exercises}
\chapter{Rational numbers, ratios, and fields}\label{ch:rationals}
\begin{readergoal}Construct fractions and learn why units, denominators, and field axioms matter in quantitative science.\end{readergoal}
\section{Fractions as equivalence classes}
Take integer pairs $(a,b)$ with $b\ne0$. Say $(a,b)\sim(c,d)$ when $ad=bc$. The resulting equivalence class represents the fraction $a/b$. Ratios like $1/2=2/4$ are different representations of the same rational number. Addition and multiplication take their familiar forms
\[\frac ab+\frac cd=\frac{ad+bc}{bd},\qquad \frac ab\cdot\frac cd=\frac{ac}{bd}.\]
Rationals form the field $\mathbb Q$: nonzero values possess multiplicative inverses, and the field operations obey associativity, commutativity, and distributivity. Division by zero is undefined: a rate with zero exposure duration is not an extreme number but an invalid operation.

\section{What does a ratio mean physically?}
Suppose a unit emits 75 spikes over 15 seconds. Its time-averaged rate is $75/15=5$ spikes/s. The numerator and denominator have different roles and units. When a software routine returns a mean across observations, ask what each observation represents and whether every observation has the same weight. The mean of within-trial rates need not equal total count divided by total exposure if trial durations differ.

\section{Rationals are dense but incomplete}
Between any two rational numbers $a<b$, their midpoint $(a+b)/2$ is another rational. Yet the equation $x^2=2$ has no rational solution. To show this, suppose $x=p/q$ is in lowest terms and $p^2=2q^2$. Then $p$ must be even; substituting $p=2k$ shows $q$ also even, contradicting coprimality. Thus the rational numbers have gaps from the viewpoint of limits and geometry.

\begin{workedexample}[Units and normalization]
Compare $\mathrm{spikes}/\mathrm{bin}$ with $\mathrm{spikes}/\mathrm{s}$. If the bin width is $\Delta t$ seconds, the rate estimate is $r=n/\Delta t$ in spikes/s. Changing $\Delta t$ changes the variance and resolution of the estimate even if the underlying experiment is unchanged. This anticipates the bias-variance trade-off in spike-rate estimation.
\end{workedexample}
\begin{exercises}
\item Explain what fails if $b=0$ is allowed in a fraction.
\item Calculate the pooled firing rate across a 5-s trial with 20 spikes and a 15-s trial with 30 spikes; compare with the unweighted average of the two trial rates.
\item Prove that between two unequal rationals there is a rational number.
\end{exercises}
\chapter{The real number system, order, and completeness}\label{ch:reals}
\begin{readergoal}Understand why the real line is more than rational arithmetic, and why continuity and convergence depend on completeness.\end{readergoal}
\section{The real line as an ordered field}
We wish to represent distances, times, and voltages without rational-number gaps. The real numbers $\mathbb R$ are an \emph{ordered field} satisfying a further completeness property. One characterization is the \emph{least-upper-bound property}: every nonempty subset of $\mathbb R$ bounded above has a least upper bound. If $A=\{x\in\mathbb R:x^2<2,\ x\ge0\}$, then $\sup A=\sqrt2$. This value cannot be supplied by rational numbers alone.

A real number may be developed rigorously using Dedekind cuts or equivalence classes of Cauchy sequences of rational numbers; either construction requires proving that the resulting operations and order are well defined. We will use the completeness axiom without adopting one construction as primary. For a fuller foundational treatment, consult \cite{MendelsonNumbers}.

\section{Bounds, supremum, and infimum}
A number $M$ is an upper bound for $A$ if $x\le M$ for every $x\in A$. The supremum is the smallest such bound; it need not belong to the set. For the open interval $(0,1)$, the supremum is 1, although 1 is not a member. The infimum is defined dually. A \emph{maximum} must be a member of the set and greater than or equal to every other member; distinguish maximum from supremum.

\section{Distances and the triangle inequality}
The absolute value defines distance on $\mathbb R$ by $d(x,y)=|x-y|$. Distance is nonnegative, symmetric, zero precisely for equal points, and obeys $d(x,z)\le d(x,y)+d(y,z)$. These four properties define a \emph{metric} on a more general space. A tolerance criterion for a numerical solution is meaningful only with a specified distance, scale, and units.

\section{Interval notation and physical uncertainty}
The interval $[a,b]$ includes endpoints, whereas $(a,b)$ excludes them. An experimental measurement can be modeled as an uncertain value lying in some interval, but the interval should not be mistaken for a statistical confidence interval without an inference procedure. A confidence interval concerns the long-run coverage of an estimator under a model; it is not simply an arbitrary range surrounding a measurement.

\begin{workedexample}[The value of completeness]
Take a sequence of rational approximations $1,1.4,1.41,1.414,\ldots$ to $\sqrt2$. Each term can be rational while its limit is not. A space of numbers meant to support convergence of such sequences must contain the limiting value. The same idea underlies convergence of numerical algorithms for signal processing.
\end{workedexample}
\begin{exercises}
\item Find the supremum, infimum, minimum, and maximum, when defined, for $A=(0,1]$.
\item Prove the reverse triangle inequality $\big||x|-|y|\big|\le |x-y|$.
\item Explain why a bounded measurement record does not imply its generating stochastic process has bounded support.
\end{exercises}
\chapter{Complex numbers: amplitude and phase}\label{ch:complex}
\begin{readergoal}Build the complex plane from real coordinates and understand Euler's formula before meeting Fourier and Hilbert transforms.\end{readergoal}
\section{Pairs with new arithmetic}
A complex number $z=a+ib$ has real part $a$ and imaginary part $b$, where $i^2=-1$. We may regard $z$ as the ordered pair $(a,b)$ of reals, with addition $(a,b)+(c,d)=(a+c,b+d)$ and multiplication $(a,b)(c,d)=(ac-bd,ad+bc)$. The conjugate is $\bar z=a-ib$ and $|z|=\sqrt{z\bar z}=\sqrt{a^2+b^2}$. For $z\ne0$, $z^{-1}=\bar z/|z|^2$.

\section{A plane with angles}
Write $z=r(\cos\theta+i\sin\theta)$ for modulus $r\ge0$ and argument $\theta$. Angles differing by integer multiples of $2\pi$ represent the same direction. Euler's formula $e^{i\theta}=\cos\theta+i\sin\theta$ gives the polar representation $z=re^{i\theta}$. Multiplication adds arguments and multiplies moduli:
\[(r_1e^{i\theta_1})(r_2e^{i\theta_2})=r_1r_2e^{i(\theta_1+\theta_2)}.\]
The periodicity of phase will require circular, rather than ordinary linear, statistics.

\section{Why complex numbers are indispensable for signals}
A sine wave can be written $x(t)=A\cos(2\pi ft+\phi)=\Re\{Ae^{i(2\pi ft+\phi)}\}$. Multiplication by $e^{i\omega t}$ shifts the angular coordinate at a constant rate. The Fourier transform treats a signal as a superposition of such rotating complex exponentials. The complex coefficient encodes two quantities: magnitude and phase. Discarding one generally loses information about the waveform.

\begin{workedexample}[Phase wraps]
The angles $179^\circ$ and $-179^\circ$ differ numerically by $358^\circ$ but lie just $2^\circ$ apart on a circle. Their arithmetic mean is $0^\circ$, which points in the opposite direction. A circular mean instead uses the argument of $e^{i179^\circ}+e^{-i179^\circ}$ and correctly points near $180^\circ$.
\end{workedexample}
\begin{exercises}
\item Compute $(1+i)^2$ and $|1+i|$.
\item Show $|zw|=|z||w|$ using complex conjugation.
\item Express $\cos\theta\cos\phi$ as a combination of complex exponentials and simplify.
\item Explain why phase averages need not be meaningful when oscillatory amplitude is close to zero.
\end{exercises}
\part{Calculus and analysis}
\chapter{Sequences, limits, and continuity}\label{ch:limits}
\begin{readergoal}Understand the precision behind ``approaches'', ``converges'', and ``continuous'' using explicit tolerances.\end{readergoal}
\section{An infinite list is a function}
The sequence $x_n=1/n$ approaches zero as $n$ grows. But what should ``approaches'' mean? We say $x_n\to L$ if, for every tolerance $\varepsilon>0$, there exists an index $N_\varepsilon$ such that $|x_n-L|<\varepsilon$ whenever $n\ge N_\varepsilon$. The order of the quantifiers matters. We must succeed for \emph{every} tolerance, although the required index may depend on the chosen tolerance.

For $x_n=1/n$, choose $N_\varepsilon>1/\varepsilon$. Then $n\ge N_\varepsilon$ implies $|1/n|<\varepsilon$. Convergence is not the same as attaining the limit; no term in this sequence equals zero. Nor is convergence established by drawing the first few points. An infinite claim requires an argument that controls all sufficiently late terms.

\section{Cauchy sequences and completeness}
A sequence is \emph{Cauchy} if its sufficiently late terms become arbitrarily close to one another: $\forall\varepsilon>0$ there exists $N$ such that $|x_m-x_n|<\varepsilon$ for all $m,n\ge N$. In a complete metric space every Cauchy sequence converges to a point of the space. The rational approximations to $\sqrt2$ form a Cauchy sequence of rationals but fail to converge \emph{within} the rationals. This shows why completeness is essential in analysis.

\section{Functions of a real variable}
We say $f(x)$ tends to $L$ as $x\to a$ if for every $\varepsilon>0$ we can find $\delta>0$ such that $0<|x-a|<\delta$ implies $|f(x)-L|<\varepsilon$. A function is continuous at $a$ when its limit equals its value at that point. Uniform continuity strengthens the statement by demanding one $\delta$ for a given $\varepsilon$ that works across the entire domain.

\section{Why this matters for continuous-time recordings}
A continuous-time model $x(t)$ does not mean its actual neuronal source changes smoothly; it means the mathematical function has a value for each real-valued time in its domain. A sampled waveform records only certain time points. Interpolating between them requires assumptions about smoothness or bandwidth. Similarly, increasing the sampling frequency is not a mathematical proof that a measurement device approximates the underlying field arbitrarily well: instrument bandwidth and noise also matter.

\begin{workedexample}[Finite-difference derivatives]
For a smooth $x(t)$, the quantity $[x(t+h)-x(t)]/h$ may approach a derivative as $h\to0$. In recorded data, however, taking ever smaller $h$ may amplify measurement noise. Convergence of an ideal mathematical expression and consistency of a noisy numerical estimator are separate questions.
\end{workedexample}
\begin{exercises}
\item Prove $2/n\to0$ from the definition.
\item Give a sequence that is bounded but does not converge.
\item Determine whether the step function $f(x)=0$ for $x<0$, $f(x)=1$ for $x\ge0$ is continuous at zero.
\item Explain the conceptual difference between a limit as $n\to\infty$ and a sample size $n$ that remains finite.
\end{exercises}
\chapter{Derivatives, integrals, and differential equations}\label{ch:calculus}
\begin{readergoal}Build rate of change and accumulation from limits and relate them to state variables and observed neural signals.\end{readergoal}
\section{A derivative is a limit of ratios}
For a differentiable function $f$, the derivative at $t$ is
\[f'(t)=\lim_{h\to0}\frac{f(t+h)-f(t)}h.\]
Its units are output units divided by input units: differentiating voltage with respect to time gives volts per second. A derivative expresses an instantaneous rate of change under a continuous model; a finite difference merely estimates this rate using discrete observations. The distinction becomes serious for noisy measurements.

\section{Integration as accumulation}
A Riemann integral $\int_a^b f(t)\,dt$ is a limit of sums of function values multiplied by small interval widths. If $r(t)$ is a firing intensity in spikes/s, its integral over $[a,b]$ has units of expected spikes, and under a conditional-intensity point-process model it represents integrated intensity, not the realized count. If $p(t)$ is power in an appropriate density scale, integrating it over time yields an accumulated quantity with corresponding units.

The fundamental theorem of calculus connects these operations: if $F(t)=\int_a^t f(s)\,ds$ and $f$ is continuous, then $F'(t)=f(t)$. Thus differentiation and integration are linked through limiting constructions.

\section{Differential equations: rules for evolving states}
A differential equation specifies how a variable changes, not necessarily its value at every time. Consider
\[\tau\frac{dx}{dt}=-x+u(t).\]
With $u(t)=u_0$ constant, the solution is $x(t)=u_0+[x(0)-u_0]e^{-t/\tau}$. The time constant $\tau>0$ controls relaxation. Without an initial condition, many trajectories satisfy the same differential equation. This is a model for first-order filtering and a building block for membrane and firing-rate dynamics.

\section{A distinction between models and data}
Recorded spikes are events and need not be represented by a differentiable function. The firing-rate function is an inferred or modeled quantity. Writing $\dot r(t)$ for the derivative of an estimated rate assumes a degree of smoothness that the raw point process does not possess. We will state these choices explicitly when we develop state-space models.

\begin{exercises}
\item Differentiate $f(t)=e^{-t/\tau}$ and confirm the result has units of inverse time.
\item Solve $dx/dt=-ax$ for $a>0$ and a specified initial condition.
\item Explain why a spike-count cumulative function is not ordinarily differentiable at spike times.
\item Compute $\int_0^T 3\,dt$ and interpret it for a rate of 3 spikes/s.
\end{exercises}
\chapter{Several variables, optimization, and geometry}\label{ch:multi}
\begin{readergoal}Generalize derivatives to gradients and Jacobians and understand curvature, critical points, and constrained optimization.\end{readergoal}
\section{The gradient is a local direction of change}
For $f:\mathbb R^n\to\mathbb R$, the gradient is $\nabla f=(\partial f/\partial x_1,\ldots,\partial f/\partial x_n)^\top$. Its dot product with a direction $v$ gives the directional derivative $D_vf=\nabla f\cdot v$ for a differentiable function. In a neighborhood of $x$, $f(x+h)\approx f(x)+\nabla f(x)^\top h$. The gradient depends on the units and coordinate scaling of the input variables.

A vector-valued map $F:\mathbb R^n\to\mathbb R^m$ has Jacobian $J_{ij}=\partial F_i/\partial x_j$. Near a point, $F(x+h)\approx F(x)+J_F(x)h$. The chain rule for $G\circ F$ is $J_{G\circ F}(x)=J_G(F(x))J_F(x)$, provided the dimensions match.

\section{Second derivatives and the Hessian}
The Hessian of a twice differentiable scalar function is the matrix of second partial derivatives $H_{ij}=\partial^2f/(\partial x_i\partial x_j)$. At a stationary point $\nabla f=0$, a positive-definite Hessian gives a strict local minimum; an indefinite Hessian indicates a saddle. The quadratic approximation
\[f(x+h)\approx f(x)+\nabla f(x)^\top h+\tfrac12 h^\top H_f(x)h\]
provides a local account of curvature and motivates Newton-type optimization.

\section{Optimization as fitting a model}
Least squares seeks $\hat\theta=\arg\min_\theta\sum_i[y_i-f_\theta(x_i)]^2$. The objective is a chosen measure of mismatch, not a guarantee that the parameter has a physiological meaning. Regularization changes the optimization problem, e.g.
\[\hat\theta=\arg\min_\theta\left\{\|y-X\theta\|_2^2+\lambda\|\theta\|_2^2\right\}.\]
The penalty limits parameter magnitude and can improve stability, but also introduces bias. Constraints such as $\theta\ge0$ are meaningful only if supported by the model and units.

\section{The manifold idea starts locally}
A smooth curve in $\mathbb R^3$ may be one-dimensional although it occupies a three-dimensional ambient space. Near a regular point, its variation can be approximated by a tangent direction. General smooth manifolds are spaces locally modeled on $\mathbb R^d$; the ambient dimension and intrinsic dimension need not coincide. PCA identifies a linear subspace, but whether the data genuinely concentrate around a nonlinear manifold is an additional question.

\begin{exercises}
\item Compute the gradient and Hessian of $f(x,y)=x^2+2y^2$.
\item Use the chain rule to differentiate $g(t)=\sin(x(t)^2)$.
\item Compare L2 and L1 penalties conceptually in regression.
\item Give an example of a curved one-dimensional set in $\mathbb R^2$.
\end{exercises}
\chapter{Integration, function spaces, and transforms}\label{ch:functionspaces}
\begin{readergoal}Prepare the mathematical ground for Hilbert spaces, Fourier coefficients, energy, and stochastic-process covariance.\end{readergoal}
\section{From vectors of numbers to functions}
The Euclidean inner product $\sum_{i=1}^n x_iy_i$ measures agreement between two finite vectors. For sufficiently integrable real functions on an interval, the analogous quantity is
\[\langle f,g\rangle=\int_a^b f(t)g(t)\,dt.\]
For complex functions we conjugate the second or first argument according to convention; here we use $\langle f,g\rangle=\int_a^b f(t)\overline{g(t)}\,dt$, linear in the first argument. The induced norm is $\|f\|_2=(\int |f|^2)^{1/2}$. Signals with finite squared integral belong to the space $L^2$, after identifying functions equal almost everywhere.

\section{Orthogonality and approximation}
Two functions are orthogonal when their inner product is zero. If $e_1,\ldots,e_k$ are orthonormal, the projection of $f$ onto their span is $\sum_j\langle f,e_j\rangle e_j$. The error is orthogonal to the span: this is the same geometry later used in least squares, PCA, and finite-basis spectral expansions.

\section{Measure and probability: what does ``almost everywhere'' mean?}
A measure assigns nonnegative sizes to suitable sets. In the familiar one-dimensional setting, interval length is the basic example. The Lebesgue integral extends Riemann integration and provides a stable theory for limits of integrable functions. ``Almost everywhere'' means except on a set of measure zero. Two functions that differ at finitely many time points can represent the same $L^2$ element, even though pointwise spike events at those instants could be scientifically significant. This reminds us that the choice of mathematical object must fit the observation type.

\section{Generalized signals and impulses}
The ideal Dirac delta $\delta(t-t_0)$ is not an ordinary $L^2$ function. In generalized-function notation its defining property is $\int f(t)\delta(t-t_0)\,dt=f(t_0)$ for suitable test functions. A spike train can be represented formally as $s(t)=\sum_j\delta(t-t_j)$, but practical recordings contain finite-time waveforms and discrete timestamps. Both descriptions are useful, and neither should be confused with the other.

\begin{workedexample}[Projection onto one sinusoid]
On $[0,T]$, let $e(t)=\sqrt{2/T}\cos(2\pi kt/T)$ for a nonzero integer $k$. Its squared integral is one. The coefficient $a=\int_0^T x(t)e(t)\,dt$ determines the best approximation $ae(t)$ to $x(t)$ within this one-dimensional subspace under squared-integral error. This is the geometric core of Fourier analysis.
\end{workedexample}
\begin{exercises}
\item Verify that $\int_0^T\sin(2\pi nt/T)\cos(2\pi mt/T)\,dt=0$ for positive integers $m,n$.
\item Explain why delta functions require care in ordinary function spaces.
\item State the difference between orthogonality and statistical independence.
\end{exercises}
\part{Linear algebra and numerical models}
\chapter{Vectors, matrices, and linear mappings}\label{ch:linear}
\begin{readergoal}Learn the grammar of population matrices before using covariance, singular values, or latent factors.\end{readergoal}
\section{Arrays are not merely spreadsheets}
A vector is an element of a vector space: an object that can be added to another vector and scaled while obeying specific algebraic rules. A matrix $A\in\mathbb R^{m\times n}$ represents a linear mapping $A:\mathbb R^n\to\mathbb R^m$. Its $j$th column is the image of the $j$th standard basis vector. The product $Ax$ forms linear combinations of these columns. Dimensions are part of the meaning: $AB$ is defined only when the number of columns of $A$ equals the number of rows of $B$.

For a population of $N$ neurons sampled at $T$ time points, set $X\in\mathbb R^{T\times N}$, with rows indexing time and columns indexing neurons. Then $Xv$ is a time series obtained by weighting neurons with $v\in\mathbb R^N$. If a recording is stored in channel-by-time orientation, as often happens in MATLAB, it must be transposed before blindly applying a formula whose convention is time-by-neuron.

\section{Rank, kernels, and invertibility}
The rank of $A$ is the dimension of its image or column space. Its kernel $\ker A=\{x:Ax=0\}$ contains directions invisible to the mapping. Rank-nullity says $\dim\ker A+\rank A=n$. A square matrix is invertible precisely when its kernel contains only zero. A measurement operator with a nontrivial kernel cannot uniquely reconstruct arbitrary inputs: different sources differ by kernel vectors while producing identical observations.

\section{Inner products and projections}
The Euclidean inner product is $\langle x,y\rangle=x^\top y$, with length $\|x\|_2=\sqrt{x^\top x}$. For a matrix $U$ whose columns are orthonormal, $U^\top U=I$, the orthogonal projection onto their span is $P=UU^\top$. Check $P^2=P$ and $P^\top=P$. The residual $(I-P)x$ is orthogonal to the selected subspace.

\section{Least squares as orthogonal projection}
If $X$ has full column rank, minimizing $\|y-X\beta\|_2^2$ yields the normal equations $X^\top X\hat\beta=X^\top y$ and solution $\hat\beta=(X^\top X)^{-1}X^\top y$. In practice, QR or SVD is usually numerically safer than explicitly forming the inverse of $X^\top X$; squaring the condition number can magnify numerical problems.

\begin{workedexample}[Which axis contains the data?]
A matrix of 30 channels and 1,000,000 samples has shape $30\times10^6$. An operation computing covariances among channels should treat the 30 channels as variables, not the million sample positions. A mistaken transpose can yield a million-by-million covariance matrix or an out-of-memory failure. Mathematical notation doubles as a memory-design contract.
\end{workedexample}
\begin{exercises}
\item Verify the projection identities for a unit vector $u$ and $P=uu^\top$.
\item Find the kernel of $A=\begin{bmatrix}1&1\\2&2\end{bmatrix}$.
\item Derive the normal equations by differentiating the squared residual norm.
\item Explain why rank deficiency compromises parameter uniqueness.
\end{exercises}
\chapter{Eigenvalues, singular values, and principal components}\label{ch:pca}
\begin{readergoal}Understand the spectrum of a matrix, its geometry, and the exact optimization solved by PCA.\end{readergoal}
\section{Eigenvectors as preserved directions}
For a square matrix $A$, an eigenvector $v\ne0$ satisfies $Av=\lambda v$. The scalar $\lambda$ describes the action on that direction. A real symmetric matrix admits an orthonormal basis of real eigenvectors and a factorization $A=Q\Lambda Q^\top$. A covariance matrix is symmetric positive semidefinite, so its eigenvalues are nonnegative.

\section{Singular value decomposition}
Every real matrix $X\in\mathbb R^{T\times N}$ admits an SVD $X=U\Sigma V^\top$ with nonnegative singular values in $\Sigma$ and orthonormal singular vectors. This construction exists even for rectangular or rank-deficient matrices. The right singular vectors describe directions in variable space; the left singular vectors describe directions in observation space. Truncating to the first $k$ components gives the best rank-$k$ approximation in Frobenius norm (Eckart--Young theorem), provided singular values are ordered descending.

\section{Deriving PCA}
Let $X_c$ denote the column-centered observations and $S=X_c^\top X_c/(T-1)$. A projected observation onto a unit vector $v$ has sample variance $v^\top Sv$. The maximum-variance direction solves
\[v_1=\arg\max_{\|v\|=1}v^\top Sv.\]
Using the Lagrangian $\mathcal L(v,\lambda)=v^\top Sv-\lambda(v^\top v-1)$ gives $Sv=\lambda v$. The maximizing direction is an eigenvector of the largest eigenvalue. Subsequent components maximize variance subject to orthogonality. If $X_c=U\Sigma V^\top$, then $\lambda_j=\sigma_j^2/(T-1)$.

\section{What PCA does not promise}
PCA maximizes sample variance, not class separation, behavioral relevance, temporal predictability, or a causal mechanism. Large variance can come from artifacts, gain fluctuations, or unmodeled experimental factors. Cross-validation, null models, and careful centering are essential. As \cite{CunninghamYu2014,Churchland2012} discuss, structured population dynamics may require objectives different from variance maximization.

\begin{workedexample}[An effective dimension]
If $\lambda_1,\ldots,\lambda_N$ are nonnegative variances, the participation ratio $d_{\rm PR}=(\sum_i\lambda_i)^2/\sum_i\lambda_i^2$ equals $N$ when all eigenvalues are equal and equals one when only one is positive. It is a continuous summary of spectral concentration, not an integer count of ``real'' sources.\end{workedexample}
\begin{exercises}
\item Derive $\lambda_j=\sigma_j^2/(T-1)$ for centered observations.
\item Compare participation ratios for $(10,0,0)$ and $(4,3,3)$.
\item Explain why replacing trials by a trial average can remove single-trial fluctuations.
\item Give a case where the leading principal component tracks recording artifact rather than neural activity.
\end{exercises}
\chapter{Tensors, reduced-rank models, and inverse problems}\label{ch:inverse}
\begin{readergoal}Move from matrix factorization to multivariate measurements, and learn how identifiable models differ from underdetermined ones.\end{readergoal}
\section{More than two indices}
An experiment with neurons $n$, time $t$, and trial $r$ may be stored as a third-order tensor $\mathcal X_{ntr}$. Its modes represent different axes of variation. A CANDECOMP/PARAFAC (CP) approximation is
\[\mathcal X_{ntr}\approx\sum_{k=1}^K a_{nk}b_{tk}c_{rk}.\]
Interpretation requires confronting scaling and permutation ambiguities: multiplying one factor by $\alpha$ and dividing another by $\alpha$ leaves the product unchanged. Additional uniqueness conditions depend on the factors and rank. The mathematical pattern is broader than one physiological explanation.

\section{Reduced-rank regression}
For predictor matrix $X\in\mathbb R^{T\times p}$ and response matrix $Y\in\mathbb R^{T\times q}$, ordinary multivariate regression fits $Y\approx XB$. Reduced-rank regression requires $\rank(B)\le r$, controlling the dimension of predictive communication. Calling the resulting dimensions ``communication subspaces'' is an interpretation built on the predictive model, not proof that specific axons mediate the influence.

\section{The forward model is not the inverse model}
Let measurements obey $y=Hx+\varepsilon$, where $x$ is a source, $H$ a measurement operator, and $\varepsilon$ noise. The \emph{forward problem} computes $y$ from a chosen source and measurement physics. The \emph{inverse problem} estimates $x$ from $y$. If $H$ has a nontrivial kernel, $x$ and $x+z$ generate the same noiseless measurement whenever $Hz=0$. This is an identifiability failure, not a limitation remedied merely by better numerical precision.

For extracellular electric potentials, source geometry and volume conduction affect $H$, and several current patterns may explain the same electrode measurements. Biophysical treatments emphasize the interpretation of forward models and the assumptions behind inverse methods \cite{Buzsaki2012,Einevoll2013}. Regularization can choose one plausible solution, but it cannot magically restore missing independent information.

\section{Regularized reconstruction}
One possible estimator is $\hat x=\arg\min_x\{\|y-Hx\|_2^2+\lambda\|Lx\|_2^2\}$. If the normal equations yield an invertible matrix, $\hat x=(H^\top H+\lambda L^\top L)^{-1}H^\top y$. The operator $L$ expresses a prior preference (e.g., smoothness), while $\lambda$ sets its influence. The result answers a conditional question: which sources fit the observations well while obeying the chosen regularity assumption?

\begin{exercises}
\item Exhibit two distinct $x\in\mathbb R^2$ satisfying $[1\ 1]x=2$.
\item Explain why a CP decomposition's factors may be arbitrarily rescaled.
\item Describe what you would need to claim that a reduced-rank relationship reflects direct anatomical connectivity.
\end{exercises}
\part{Probability, statistics, and time}
\chapter{Probability begins with sets of possible outcomes}\label{ch:probability}
\begin{readergoal}Translate uncertainty into an explicit sample space, events, conditional probabilities, and independence assumptions.\end{readergoal}
\section{Experiments and events}
A random experiment has possible outcomes in a \emph{sample space} $\Omega$. An \emph{event} is a subset $A\subseteq\Omega$. A probability measure $P$ assigns numbers to events with $P(A)\ge0$, $P(\Omega)=1$, and countable additivity across disjoint events. This connects the set operations learned in Chapter~\ref{ch:sets} to the calculus of uncertainty. A model of stochastic spike counts has a different sample space from a model of electrode voltage traces; neither should be silently substituted for the other.

\section{Conditional probability}
When $P(B)>0$, define $P(A\mid B)=P(A\cap B)/P(B)$. Rearranging gives the multiplication rule $P(A\cap B)=P(A\mid B)P(B)$. If $A$ and $B$ are independent, $P(A\cap B)=P(A)P(B)$. Independence is a model property, not merely a statement that the observed sample correlation vanishes.

The law of total probability for a partition $(B_j)$ is $P(A)=\sum_jP(A\mid B_j)P(B_j)$. Bayes' theorem follows:
\[P(B_j\mid A)=\frac{P(A\mid B_j)P(B_j)}{\sum_kP(A\mid B_k)P(B_k)}.\]
In Bayesian analysis, a prior over model parameters and a likelihood for observed data yield a posterior. The posterior is conditional on both the likelihood and the prior; a credible interval and a frequentist confidence interval do not have interchangeable interpretations.

\section{Conditioning changes the question}
Two neurons can appear correlated because each varies with movement speed. Conditioning on behavior may reduce this apparent association. But conditioning can also introduce new dependencies, as in collider bias. A scientific claim must specify what was conditioned on, why, and whether the variables were measured accurately. The conditional modeling framework for spike jitter makes this especially explicit \cite{Amarasingham2012}.

\begin{workedexample}[A base-rate effect]
Suppose only 1 percent of candidate events are genuine. A detector identifies 90 percent of genuine events and falsely flags 5 percent of nongenuine events. The probability that a flagged event is genuine is $0.90(0.01)/[0.90(0.01)+0.05(0.99)]\approx0.154$. High sensitivity does not by itself imply high positive predictive value.
\end{workedexample}
\begin{exercises}
\item Derive $P(A\cup B)=P(A)+P(B)-P(A\cap B)$.
\item Give independent events and dependent events from a six-sided die experiment.
\item Explain why a detector's posterior event probability depends on prevalence.
\end{exercises}
\chapter{Random variables, expectation, and statistical moments}\label{ch:moments}
\begin{readergoal}Define distributions and calculate moments before discussing firing-rate variability or moving-window envelope statistics.\end{readergoal}
\section{A random variable is a function}
A random variable $X:\Omega\to\mathbb R$ assigns a numerical value to each possible outcome, subject to a measurability requirement that ensures events such as $\{\omega:X(\omega)\le x\}$ have defined probabilities. Its cumulative distribution function is $F_X(x)=P(X\le x)$. Discrete random variables have probability mass functions; sufficiently regular continuous variables have probability densities. A density is not itself a probability; probabilities arise through integration over regions.

\section{Means, variance, skewness, kurtosis}
The mean is $\mu=E[X]$; variance is $\sigma^2=E[(X-\mu)^2]$; standardized skewness is $\gamma_1=E[(X-\mu)^3]/\sigma^3$; and standardized kurtosis is $\gamma_2=E[(X-\mu)^4]/\sigma^4$. Some authors report excess kurtosis $\gamma_2-3$, which equals zero for a Gaussian distribution. If $\sigma=0$, standardized higher moments are undefined. Higher moments are sensitive to rare large amplitudes and may require large samples for stable estimation.

For a time-indexed envelope $A(t)$, a sliding-window empirical mean is $\hat\mu(t)=\frac1{M}\sum_{j=0}^{M-1}A(t+j\Delta t)$ with analogous central moments. If consecutive windows overlap heavily, the resulting statistics are strongly dependent: sampling an estimate every 100 ms does not create a new independent 10-s experiment at every step.

\section{Covariance and correlation}
Covariance is $\Cov(X,Y)=E[(X-EX)(Y-EY)]$. Its scale depends on units, while Pearson correlation normalizes by standard deviations: $\rho=\Cov(X,Y)/(\sigma_X\sigma_Y)$ for nonzero variances. A zero correlation does not guarantee independence in general. Let $X$ be symmetric around zero and $Y=X^2$; then $\Cov(X,Y)=E[X^3]-E[X]E[X^2]=0$, although $Y$ is determined by $X$.

\section{The law of total variance}
For random variables $X,Z$ with finite variance,
\[\Var(X)=E[\Var(X\mid Z)]+\Var(E[X\mid Z]).\]
The first term measures variation remaining after conditioning; the second measures variation in conditional means. This identity is an essential bridge to decomposing neuronal trial-to-trial variability into point-process and gain contributions \cite{Goris2014}.

\begin{workedexample}[Poisson versus overdispersed counts]
For $N\sim\mathrm{Poisson}(\lambda)$, $E[N]=\Var(N)=\lambda$. If $N\mid\Lambda\sim\mathrm{Poisson}(\Lambda)$ and $\Lambda$ varies, then $\Var(N)=E[\Lambda]+\Var(\Lambda)$. Extra variance can arise from latent fluctuations even if spikes are conditionally Poisson.
\end{workedexample}
\begin{exercises}
\item Calculate the population mean, variance, skewness, and kurtosis of the discrete values $0,1,2$ under equal probabilities.
\item Prove $\Var(aX+b)=a^2\Var(X)$.
\item Explain the difference between window overlap and an actual temporal correlation in the source.
\item Derive the overdispersion identity from total variance.
\end{exercises}
\chapter{Inference, likelihood, and stochastic processes}\label{ch:stochastic}
\begin{readergoal}Separate population parameters from estimators, and connect time-dependent random variables to autocorrelation and stationarity.\end{readergoal}
\section{Data, models, estimators}
An unknown parameter $\theta$ describes a proposed data-generating distribution $p(x\mid\theta)$. An \emph{estimator} $\hat\theta(X_1,\ldots,X_n)$ is a function of random data; its realized value is an estimate. An estimator can be biased, have high variance, or be inconsistent. Mean squared error decomposes as $\mathrm{MSE}(\hat\theta)=\Var(\hat\theta)+[E\hat\theta-\theta]^2$ when these quantities exist. Estimation is therefore not reduced to maximizing an observed fit score.

The likelihood $L(\theta;x)=p(x\mid\theta)$ views the observed data as fixed and the parameter as variable. For independent Poisson counts $n_i$ with intensities $\lambda_i(\theta)$, the log-likelihood is
\[\ell(\theta)=\sum_i\{n_i\log\lambda_i(\theta)-\lambda_i(\theta)-\log(n_i!)\}.\]
The assumption of independent counts must be justified or modified. A large likelihood does not prove the model is correct, and a small $p$-value does not give the probability that the null hypothesis is true.

\section{Time indexed random variables}
A stochastic process is a family $(X_t)_{t\in T}$ of random variables on a common probability space. For weak stationarity, the mean is constant and the covariance between $X_t$ and $X_s$ depends only on their lag $t-s$ (assuming finite second moments). Define autocovariance $C_X(\tau)=\Cov(X_t,X_{t+\tau})$. Its Fourier transform, under appropriate conditions, gives a power spectral density. Stationarity is a property of the model, not something one can prove from one short recording.

\section{Conditional structure and hidden state}
A state-dependent process may have time-varying conditional mean or variance despite an apparently stable recording. A latent state $Z_t$ might represent arousal, behavior, or an abstract computational state. The distinction between $p(X_t\mid Z_t)$ and the marginal $p(X_t)$ explains why pooling states can create spurious correlations and overdispersion. This is a key insight in interpreting correlations \cite{CohenKohn2011,Goris2014}.

\begin{workedexample}[A stochastic oscillator is not just a sine wave]
A pure sinusoid has a well-defined frequency and deterministic phase evolution. A noisy oscillatory process may have a broad spectral peak and time-varying phase. Conversely, a stochastic process with $1/f$-like spectrum can place substantial power within the theta band without containing an independent theta oscillator. The distinction motivates explicit spectral modeling \cite{Donoghue2020,Gerster2022}.
\end{workedexample}
\begin{exercises}
\item Explain what changes when observations are serially correlated but an iid likelihood is used.
\item Derive the mean and variance of an average of $n$ independent variables with equal finite variance.
\item Define weak stationarity in your own words and give a counterexample.
\item Explain why trial averaging can suppress evidence of state switches.
\end{exercises}
\part{Mathematics of electrophysiological signals}
\chapter{From a physical voltage to a sampled sequence}\label{ch:sampling}
\begin{readergoal}Understand the measurement chain, units, sample rates, aliasing, and the epistemic difference between source and signal.\end{readergoal}
\section{What electrodes actually measure}
Extracellular electrodes record potential differences relative to a reference. Their outputs reflect contributions of multiple transmembrane currents, electrode geometry, and conductive tissue. The recorded LFP is not a direct readout of a particular neuron or synapse. A physical forward model connects currents to voltages \cite{Buzsaki2012,Einevoll2013}. Analog amplification, referencing, and filtering occur before the digitized values are stored.

\section{Sampled time is a lattice}
For sampling frequency $f_s$, write $\Delta t=1/f_s$ and $x[n]=x(n\Delta t)$ for integer $n$ in the acquisition index set. The Nyquist frequency is $f_s/2$. A sinusoid at a frequency above Nyquist can produce the same discrete samples as a lower-frequency sinusoid; this is \emph{aliasing}. A pre-sampling low-pass filter is therefore required to suppress frequencies that would fold into the band of interest. Merely applying a digital low-pass filter after sampling cannot undo aliasing already introduced.

\section{The ideal sampling theorem and its assumptions}
If a continuous signal is strictly bandlimited to frequencies $|f|<f_s/2$, its samples determine the original signal through ideal sinc interpolation. Actual brain recordings are neither exactly bandlimited nor free of instrumentation artifacts. The theorem establishes a mathematical possibility under stated conditions, not a promise that every sampled waveform reproduces all physical current sources.

\section{Storage, quantization, and memory mapping}
A discrete recording is represented in finite precision. With $C$ channels, $N$ samples, and $b$ bytes per stored sample, uncompressed size is $C N b$ bytes (ignoring metadata and headers). Memory-mapped storage can support access to selected time blocks without allocating the entire dataset. It does not guarantee fast noncontiguous access: array layout, strides, page faults, and disk bandwidth still matter. Always document sample type, endian convention, orientation, sampling rate, units, and channel metadata.

\begin{workedexample}[A large acquisition]
For 128 channels at 30,000 samples/s over one hour, stored as signed 16-bit integers, the raw size is $128\cdot30000\cdot3600\cdot2=27{,}648{,}000{,}000$ bytes (about 27.65 GB in decimal units). Converting the entire array to 64-bit floats multiplies memory by four. A streaming analysis therefore need not sacrifice mathematical precision in its intermediate accumulators, but should avoid careless full-array expansion.
\end{workedexample}
\begin{exercises}
\item Determine the Nyquist frequency for an LFP sampled at 1250 Hz.
\item Calculate the storage size of 80 channels sampled at 2 kHz for 2 hours using float32.
\item Explain why the reference electrode can affect estimated coherence.
\item Distinguish a physical source model, an amplifier signal, and a floating-point array.
\end{exercises}
\chapter{Fourier analysis, convolution, and digital filtering}\label{ch:fourier}
\begin{readergoal}Derive spectral representations and see why filtering is a temporal operation with potential artifacts.\end{readergoal}
\section{The discrete Fourier transform}
For $N$ discrete samples, define
\[X[k]=\sum_{n=0}^{N-1}x[n]e^{-i2\pi kn/N},\qquad k=0,\ldots,N-1.\]
The inverse is $x[n]=N^{-1}\sum_{k=0}^{N-1}X[k]e^{i2\pi kn/N}$. The frequency index represents $f_k=kf_s/N$ modulo $f_s$. Frequency resolution in the ordinary DFT grid is $\Delta f=f_s/N=1/T$ where $T=N/f_s$. Zero-padding interpolates the frequency grid but does not create new resolving power for nearby components.

Parseval's identity for this convention is $\sum_n|x[n]|^2=N^{-1}\sum_k|X[k]|^2$. One-sided and two-sided PSD conventions differ by factors of two away from DC and Nyquist; implementations must state normalization and units. The PSD estimates power per hertz, not simply squared FFT amplitude.

\section{Convolution and linear time-invariant systems}
A discrete linear time-invariant filter has output $y[n]=\sum_m h[m]x[n-m]$, where $h$ is its impulse response. Under appropriate conditions, $Y(e^{i\omega})=H(e^{i\omega})X(e^{i\omega})$. The transfer function has both magnitude and phase. Its impulse and step responses show how brief transients can be smeared in time or transformed into artificial-looking oscillations \cite{CheveigneNelken2019}.

\section{FIR, IIR, and zero-phase filtering}
A finite impulse-response (FIR) filter has a finite set of coefficients; an infinite impulse-response (IIR) filter includes feedback. Applying a forward-backward filter can cancel net phase shift in the interior of a record but is noncausal and changes the effective magnitude response. Boundaries require special handling; filtering a short isolated segment can produce misleading edge effects. In an online decoder, a filter that uses future samples is not available without added latency.

\section{Cross-correlation and spectral duality}
For a weakly stationary process, the autocovariance $C_{xx}[\ell]=E[(X_n-\mu)(X_{n+\ell}-\mu)]$ is related to the PSD through a Fourier transform (Wiener--Khinchin). This relationship creates a bridge between lag-domain dependencies and frequency-domain power. Oscillatory-looking autocovariance and a spectral peak are related descriptions; neither alone proves the existence of a distinct biological oscillator.

\begin{workedexample}[Moving-average low-pass filter]
Take $y[n]=(x[n]+x[n-1])/2$. Its frequency response is $H(e^{i\omega})=\tfrac12(1+e^{-i\omega})=e^{-i\omega/2}\cos(\omega/2)$. Thus magnitude attenuation and phase delay both emerge from an apparently simple average. At Nyquist ($\omega=\pi$), the response is zero.
\end{workedexample}
\begin{exercises}
\item Calculate the DFT of a constant sequence and identify its DC term.
\item Derive the moving-average frequency response for three samples.
\item Explain the difference between a Fourier coefficient and a PSD estimate.
\item Design a simulation exposing a filter-induced false oscillatory event.
\end{exercises}
\chapter{Spectrograms, multitapers, and periodic-versus-aperiodic spectra}\label{ch:spectral}
\begin{readergoal}Distinguish a spectral estimate from the generating spectrum and avoid inferring oscillations from arbitrary band power.\end{readergoal}
\section{Why spectral estimates fluctuate}
A finite realization contains sampling variability and boundary effects. The periodogram's frequency resolution changes with recording duration, and its variance does not simply vanish by taking a long record without appropriate averaging. Windowing reduces leakage at the price of altered resolution. Welch estimates average periodograms across segments; the segments may overlap, but this does not make them independent replicates.

\section{Spectrograms and the time-frequency trade-off}
The short-time Fourier transform is
\[X(\tau,f)=\int_{-\infty}^{\infty}x(t)w(t-\tau)e^{-i2\pi ft}\,dt,\]
where $w$ is a localized window. Narrow time windows localize fast transients while broadening spectral resolution. Long windows sharpen frequency information but smear event timing. A spectrogram $|X(\tau,f)|^2$ is not an instantaneous power measurement at a mathematical point: it summarizes a neighborhood determined by the window.

\section{Multitaper estimation}
The multitaper method forms spectral estimates using several mutually orthogonal tapers with controlled spectral concentration. Discrete prolate spheroidal sequences are commonly selected for a time-bandwidth product $NW$; the number of useful tapers is often on the order of $2NW-1$. The smoothing bandwidth and record length must be reported. Multitaper methods are central to the Chronux approach \cite{Bokil2010}.

\section{Spectral peaks over the background}
Brain spectra often contain aperiodic structure approximated locally by a $1/f^\beta$ law, along with putative narrowband components. A rise in band power may result from an oscillatory peak, a change in the aperiodic offset, a change in the exponent, or a frequency shift. Models such as FOOOF/specparam parameterize periodic and aperiodic components; IRASA uses irregular resampling to estimate the aperiodic part. These methods have distinct failure modes, especially when peaks are wide or overlap \cite{Donoghue2020,Gerster2022}.

\begin{workedexample}[Two indistinguishable band powers]
Imagine condition A has a 10-Hz Gaussian peak over a low background while condition B has no distinct peak but a sufficiently elevated broadband background. The integrals of their spectral power from 8 to 12 Hz may be equal. The same band-power summary corresponds to different generative accounts; it is therefore not a sufficient statistic for the existence of an oscillator.
\end{workedexample}
\begin{exercises}
\item Explain the meaning of time-bandwidth product in multitaper estimation.
\item Simulate a signal whose $1/f$ slope changes but whose oscillatory peak remains fixed.
\item Explain why a spectral peak detector needs a null or baseline definition.
\item Describe two ways spectrogram edge artifacts can mislead interpretation.
\end{exercises}
\chapter{Analytic signals, circular statistics, and cross-frequency coupling}\label{ch:hilbert}
\begin{readergoal}Connect complex phase to the analytic signal and understand exactly what phase-amplitude coupling measures.\end{readergoal}
\section{The analytic signal}
For a real-valued signal $x(t)$, define its Hilbert transform $\mathcal H\{x\}(t)$ by a principal-value convolution with $1/(\pi t)$, under suitable function-space assumptions. The analytic signal is
\[z(t)=x(t)+i\mathcal H\{x\}(t)=a(t)e^{i\phi(t)},\]
where $a(t)=|z(t)|$ is amplitude and $\phi(t)=\arg z(t)$ is phase. In Fourier space, analytic-signal construction suppresses negative frequencies and doubles positive frequencies (with appropriate DC and Nyquist conventions). These operations acquire a clean oscillatory interpretation for a reasonably narrowband component. Phase can be unstable when instantaneous amplitude is very small.

\section{Filtering comes before interpretation}
Often the signal is band-pass filtered to $x_b(t)$ before Hilbert transformation. The resulting $a_b(t)$ and $\phi_b(t)$ describe the \emph{filtered component}, not a unique property of the raw potential. Choice of passband, filter order, padding, and transitions can change the estimates. For nonsinusoidal waves, harmonics introduced by waveform shape can appear as high-frequency power linked to the low-frequency phase \cite{ColeVoytek2018}.

\section{Circular means and concentration}
For phases $\phi_1,\ldots,\phi_n$, set $R=n^{-1}\sum_j e^{i\phi_j}$. The mean direction is $\arg R$ when $R\ne0$, and the mean resultant length $|R|$ measures concentration between zero and one. Classical tests such as Rayleigh's test target nonuniformity in directions under specified assumptions. They do not automatically correct for temporal autocorrelation or repeated units from the same animal \cite{Berens2009}.

\section{Phase locking and phase-amplitude dependence}
The phase-locking value of phase differences is $\mathrm{PLV}=|n^{-1}\sum_j e^{i(\phi_{1j}-\phi_{2j})}|$. It measures concentration of phase differences, not amplitude modulation. A PAC statistic such as the mean vector length $|n^{-1}\sum_j A_{\mathrm{high}}(t_j)e^{i\phi_{\mathrm{low}}(t_j)}|$ quantifies phase dependence of a high-band amplitude; its scale and biases require normalization and surrogates. The modulation index based on divergence of binned amplitude distributions is a different statistic. These quantities should not be conflated.

\section{The interpretation trap}
A significant PAC value need not represent two physiologically distinct oscillators interacting. Common nonstationarity, sharp transients, harmonics, or filtering artifacts can generate cross-frequency relationships \cite{Aru2015}. Good practice checks waveform shape, aperiodic structure, trial boundaries, and meaningful surrogate null models. Early PAC observations in human cortex motivate the question, but they do not eliminate the general identifiability issues \cite{Canolty2006}.

\begin{workedexample}[One nonsinusoidal oscillator]
Consider the following waveform:
\[x(t)=\cos(2\pi f_0 t)+0.5\cos(4\pi f_0 t).\]
Its higher-frequency harmonic is phase-determined by the same fundamental oscillator. Filtering separately near $f_0$ and $2f_0$ produces components whose phase relationship is stable. It would be misleading to infer two interacting independent biological oscillators from this constructed example.
\end{workedexample}
\begin{exercises}
\item Show why the arithmetic mean of angles near $\pm\pi$ is unreliable.
\item Distinguish PLV from PAC with explicit input variables.
\item Propose a surrogate method and specify exactly what structure it preserves.
\item Explain why zero-phase filters cannot be interpreted as real-time causal processing.
\end{exercises}
\chapter{Cross-spectra, coherence, and Granger prediction}\label{ch:connectivity}
\begin{readergoal}Derive dependence measures for multivariate signals and understand why prediction is not intervention.\end{readergoal}
\section{From autocovariance to cross-covariance}
For weakly stationary processes $X_t$ and $Y_t$, define $C_{xy}(\tau)=E[(X_t-\mu_x)(Y_{t+\tau}-\mu_y)]$. Its Fourier transform is the cross-spectral density $S_{xy}(f)$. The autospectra $S_{xx},S_{yy}$ describe each process; the cross-spectrum is complex and contains magnitude and phase information about their lagged relationship.

Magnitude-squared coherence is
\[\mathcal C_{xy}(f)=\frac{|S_{xy}(f)|^2}{S_{xx}(f)S_{yy}(f)},\]
when the denominator is positive. For valid spectral-density matrices, $0\le\mathcal C_{xy}(f)\le1$. Coherence is non-directional. Magnitude-only summaries discard phase and cannot by themselves indicate whether one observed signal drives another.

\section{Common input and reference contamination}
Suppose $x=s+n_x$ and $y=s+n_y$ share a signal $s$ while $n_x,n_y$ are independent noise processes. Then $x$ and $y$ can exhibit substantial coherence even though no direct interaction between their unique sources exists. Similarly, a reference electrode that carries electrical activity can insert a common component into many recordings. Bipolar referencing or spatial derivatives may reduce some effects but introduce others; one must evaluate the measurement geometry and signal-to-noise ratio \cite{Bastos2016,CohenTsuchiya2018}.

\section{Autoregressive models and direction of prediction}
A vector autoregressive model of order $p$ for $z_t\in\mathbb R^d$ is
\[z_t=\sum_{k=1}^p A_kz_{t-k}+\varepsilon_t.\]
Consider predicting $Y_t$ from its own past versus predicting it from both its own past and the past of $X_t$. A classical Gaussian linear Granger statistic compares residual variances:
\[F_{X\to Y}=\log\frac{\sigma^2_{\rm restricted}}{\sigma^2_{\rm full}}.\]
When the full model predicts better, the statistic is positive. The arrow denotes an improvement in prediction under the candidate model, not proof that manipulating $X$ would change $Y$. Latent common causes, instantaneous mixing, nonstationarity, model order, and noise can alter apparent direction. Spectral Granger extensions decompose predictive relationships by frequency under additional factorization assumptions \cite{Cekic2018}.

\section{Transfer entropy and conditional dependence}
Transfer entropy is a conditional information measure, often written $I(Y_t;X_{t-1}^{(p)}\mid Y_{t-1}^{(q)})$. It compares how informative the past of $X$ is about $Y_t$ once the past of $Y$ is known. Nonlinear estimators can be useful but require much data and careful treatment of bias. Prediction-based asymmetry is not a universal causal-identification guarantee.

\begin{workedexample}[A shared driver]
Let $s_t$ be a stochastic signal, $x_t=s_t+\eta_t$, and $y_t=s_{t-1}+\xi_t$ with independent observation noises. The lag creates directional predictability from $x$ to $y$, but the actual source is $s$, not a connection from the observed $x$ channel to the observed $y$ channel. This is a constructive counterexample to an unrestricted causal interpretation.
\end{workedexample}
\begin{exercises}
\item State why coherence is unchanged by a fixed nonzero scaling of one signal.
\item Distinguish phase locking, coherence, and Granger prediction.
\item Construct an example of common reference that raises coherence without coupling.
\item Describe checks for VAR residual autocorrelation and stability.
\end{exercises}
\chapter{Extracellular potentials and source reconstruction}\label{ch:biophysics}
\begin{readergoal}Connect physical current sources, volume conduction, referencing, and the limitations of current-source-density analysis.\end{readergoal}
\section{Potential is a field}
An extracellular electric potential $\phi(\mathbf r,t)$ is a scalar field. The electric field is $\mathbf E=-\nabla\phi$. Electrode voltages are differences in potential; changing the reference affects the data. In an idealized homogeneous isotropic quasistatic conductor, a point current source contributes potential proportional to $I/(4\pi\sigma r)$, where $\sigma$ is conductivity and $r$ the distance to the source. This approximation illustrates spatial spreading but does not accurately describe every layered or anisotropic brain tissue.

\section{Current source density}
Conservation of charge and Ohm's law under a quasistatic conductor model lead to relationships between current source density and spatial derivatives of potential. Under constant scalar conductivity, the extracellular current-source-density convention can be written $\mathrm{CSD}=-\sigma\nabla^2\phi$ (the sign depends on the source/sink convention). A laminar one-dimensional approximation replaces the Laplacian by the second depth derivative; a discrete central difference is
\[\frac{\partial^2\phi}{\partial z^2}\Big|_{z_j}\approx\frac{\phi_{j+1}-2\phi_j+\phi_{j-1}}{(\Delta z)^2}.\]
This formula assumes uniform spacing and sufficiently smooth fields. Differentiation amplifies high-spatial-frequency noise, often making smoothing or regularization necessary.

\section{Assumptions should travel with the result}
The physics depends on conductivities, boundaries, geometry, electrode positions, and reference schemes. A depth reversal in LFP phase may be consistent with a laminar generator but does not alone locate a unique set of synaptic currents. Modeling and validation with synthetic sources help distinguish plausible mechanisms \cite{Buzsaki2012,Einevoll2013}.

\begin{exercises}
\item Derive the units of $\sigma\nabla^2\phi$ in SI units.
\item Explain why a second spatial derivative is sensitive to electrode noise.
\item State two conditions needed for a one-dimensional laminar CSD approximation.
\end{exercises}
\part{Point processes and neuronal populations}
\chapter{Spike trains as point processes}\label{ch:points}
\begin{readergoal}Model spike timestamps without pretending they are uniformly sampled voltages.\end{readergoal}
\section{Events rather than amplitudes}
A spike train is a locally finite collection of event times $\{t_1,\ldots,t_K\}$. Its counting process is $N(t)=\#\{j:t_j\le t\}$. It is nondecreasing, integer valued, and jumps at event times. In generalized-function notation $dN(t)/dt=\sum_j\delta(t-t_j)$, but the counting process is the more elementary construction.

\section{Poisson models and their scope}
For a homogeneous Poisson process with rate $\lambda>0$, counts in disjoint intervals are independent and $N(b)-N(a)\sim\mathrm{Poisson}(\lambda(b-a))$. The waiting times are exponentially distributed, and the memoryless model has no refractory period. Real neurons often violate some of these assumptions. A time-inhomogeneous Poisson model admits $\lambda(t)$, with integrated intensity $\int_a^b\lambda(s)ds$ over an interval.

A conditional intensity $\lambda(t\mid\mathcal H_t)$ is a limiting event rate given the observed history $\mathcal H_t$. Roughly,
\[\lambda(t\mid\mathcal H_t)=\lim_{\Delta\downarrow0}\frac{P(N(t+\Delta)-N(t)=1\mid\mathcal H_t)}{\Delta},\]
when the limit is defined and multiple jumps are negligible at small scales. History dependence permits refractory behavior, bursting, and coupling covariates.

\section{Binning and the loss of exact timing}
Dividing time into bins of width $\Delta$ creates counts $n_k=N((k+1)\Delta)-N(k\Delta)$. This may simplify statistics but loses information about positions inside bins. Narrower bins increase temporal precision yet yield sparse counts; wider bins stabilize rate estimates but obscure precise coincidence. A precise synchrony claim therefore needs a time-resolution convention and an explicit null model.

\begin{workedexample}[Poisson probability of a silent window]
A 10-spikes/s homogeneous Poisson process has expected count $\lambda\Delta=0.1$ in 10 ms. The chance of observing zero spikes is $e^{-0.1}\approx0.905$. Frequent silent bins are compatible with a nonzero underlying intensity.
\end{workedexample}
\begin{exercises}
\item Derive the probability of at least one spike in a window under a homogeneous Poisson model.
\item Explain why spike sorting uncertainty can affect autocorrelograms.
\item Distinguish an observed spike count from conditional intensity.
\item Give a point-process model that excludes spikes for a fixed period after each event.
\end{exercises}
\chapter{Spike-train GLMs, resampling, and conditional hypotheses}\label{ch:glm}
\begin{readergoal}Build spike prediction models and valid null tests from conditional probabilities instead of selecting statistics first.\end{readergoal}
\section{A Poisson generalized linear model}
A GLM relates the expected count to covariates through a link function. For binned counts $n_k$ over exposure $\Delta_k$, one specification is
\[n_k\mid x_k\sim\mathrm{Poisson}(\mu_k),\qquad\log\mu_k=\log\Delta_k+\beta_0+x_k^\top\beta.\]
The offset $\log\Delta_k$ ensures that changing bin exposure has the correct interpretation in rate units. Under this model, $\lambda_k=\exp(\beta_0+x_k^\top\beta)$ is the conditional event rate. For sufficiently fine bins a Bernoulli model with appropriate link can also be used, but the likelihood assumptions differ.

\section{History filters and coupling}
A point-process GLM can augment predictors with spike-history bases or other units' recent spikes. Negative recent-self-history weights may express refractoriness; positive interactions may express predictive coupling. The estimated coefficients depend on basis choice, observation artifacts, unobserved common inputs, and regularization. Coupling filters are model parameters and not direct measurements of synaptic transmission.

\section{Model assessment rather than fit alone}
Inspect residual or time-rescaled-event structure, held-out log-likelihood, calibration, stability across bin widths, and sensitivity to covariates. Time-rescaling transforms integrated intensities between consecutive events; a correctly specified conditional-intensity model yields rescaled intervals with a known reference distribution. Failure suggests misspecification but does not identify the correct alternative automatically.

\section{Jitter and conditional resampling}
A resampling null model may fix the spike count inside each jitter window and randomize event locations within specified allowed regions. Exactness of conditional tests depends on the null distribution used for resampling. Generic ``jitter the spikes'' is not a statistical procedure until its preserved quantities, randomized quantities, boundary handling, and sampling distribution are fully specified. Conditional modeling clarifies the hypotheses supported by interval jitter and related tests \cite{Amarasingham2012}.

If $T_{\rm obs}$ is a test statistic and $T_1^*,\ldots,T_B^*$ are exchangeable null replicates under the test construction, a common Monte Carlo tail estimate is
\[\hat p=\frac{1+\sum_{b=1}^B\mathbf1\{T_b^*\ge T_{\rm obs}\}}{B+1}.\]
The plus-one convention avoids zero estimated $p$-values, but validity also requires the simulation to represent the intended null.

\begin{workedexample}[What is preserved?]
Suppose one window has three spikes. A conditional jitter method preserving this count cannot test whether the process increased its number of spikes in that window; it can test aspects of fine timing within the window. Rejection must be interpreted relative to preserved coarse structure, not as evidence for a particular synapse.
\end{workedexample}
\begin{exercises}
\item Derive the score of the Poisson log-likelihood with respect to $\beta$.
\item State the null hypothesis and exact conditioning in a 20-ms interval jitter test.
\item Explain why randomizing every timestamp independently across a session may destroy behavioral modulation.
\item Design a split-trial validation plan that avoids training and testing on overlapping windows from the same trial.
\end{exercises}
\chapter{Neuronal covariance and the geometry of population responses}\label{ch:population}
\begin{readergoal}Distinguish signal correlation, noise correlation, factor structure, and the information represented in population covariance.\end{readergoal}
\section{One neuron versus a population}
Let $r(t)=(r_1(t),\ldots,r_N(t))^\top$ represent firing-rate estimates for $N$ units. A time series $r(t)$ traces a curve through $\mathbb R^N$. Neurons' responses may change together because of shared inputs, latent states, task drive, or interactions. The covariance matrix captures second-order cofluctuation but does not identify which explanation generated it.

In repeated-stimulus experiments, \emph{signal correlation} measures similarity in mean tuning across stimuli; \emph{noise correlation} measures trial-to-trial cofluctuation after controlling the stimulus or condition. These are distinct statistics, and their estimates depend on temporal window, firing rate, spike sorting, and brain state \cite{CohenKohn2011}.

\section{Shared gain and overdispersion}
A modulated Poisson count model can be written $N\mid G\sim\mathrm{Poisson}(G\mu)$, where $G>0$ is a random gain with $E[G]=1$. Then $E[N]=\mu$ and
\[\Var(N)=\mu+\mu^2\Var(G).\]
The second term represents fluctuations in the latent gain; the decomposition follows directly from the law of total variance. For two neurons with gains $G_1,G_2$, shared gain fluctuations can generate positive count covariance. This does not require a direct synaptic link between their spike generators \cite{Goris2014}.

\section{Latent factors and factor analysis}
A linear observation model $r_t=Cz_t+d+\varepsilon_t$ represents high-dimensional responses through lower-dimensional latent $z_t$. Factor analysis separates low-rank shared covariance from modeled observation noise. GPFA adds smooth temporal covariance priors to latent trajectories. Identifying the effective latent dimensionality requires validation and attention to scaling, stationarity, and neuron subsampling \cite{CunninghamYu2014,Gao2017}.

\section{Random-matrix caution}
For a high-dimensional sample covariance, random eigenvalue fluctuations can be large even when the true covariance is isotropic. In a simple independent-Gaussian null with feature-to-sample ratio $q$, the Mar\v{c}enko--Pastur bulk has support $\sigma^2(1\pm\sqrt q)^2$ under the standard asymptotic assumptions. This is not a universal biological noise threshold; correlated trials, heterogeneous variances, and temporal smoothing violate simple null assumptions.

\begin{exercises}
\item Derive the gain-modulated variance formula using total variance.
\item Explain how simultaneous recording differs from assembling pseudo-populations across sessions.
\item Give a counterexample where noise correlation increases without a change in the direct connection between neurons.
\item Explain why a large first eigenvalue is not automatically a genuine latent neural mode.
\end{exercises}
\chapter{State-space dynamics, neural manifolds, and rotations}\label{ch:dynamics}
\begin{readergoal}Connect dynamical evolution, latent inference, and geometric structure while keeping descriptive and mechanistic claims separate.\end{readergoal}
\section{A state is more than a snapshot}
A discrete linear dynamical system has $z_{t+1}=Az_t+Bu_t+w_t$ and observations $y_t=Cz_t+v_t$. The hidden state $z_t$ evolves according to a transition rule while observations mix it with noise. If the model is linear Gaussian, Kalman filtering estimates the posterior state sequentially. Smoothing uses future as well as past observations and is not equivalent to an online estimate.

\section{Continuous dynamics and stability}
A continuous linear system obeys $\dot z=Az+Bu(t)$. The eigenvalues of $A$ govern local behavior: negative real parts give asymptotic decay for unforced linear systems, while complex eigenvalues generate damped rotational components. More generally $\dot z=f(z,u,t)$ can represent nonlinear dynamics. Stable equilibria and attractors are properties of the chosen model; inferring them from noisy, partial observations requires care.

\section{jPCA and rotational population trajectories}
Principal-component analysis selects directions of large variance, not necessarily directions displaying dynamics. The jPCA approach first reduces dimensionality, then fits an approximately skew-symmetric linear flow for the derivatives of the reduced trajectories, highlighting rotations \cite{Churchland2012}. The existence of a rotational projection is a description of population activity under chosen preprocessing; attributing that rotation to a particular recurrent mechanism requires additional evidence.

\section{Neural manifolds and coordinate nonuniqueness}
A low-dimensional manifold is a set locally parameterized by a smaller number of coordinates than its ambient dimension. If observations obey $y=g(z)$ with $z\in\mathbb R^d$, many coordinate systems for $z$ can describe the same observed trajectories. Thus latent axes may not have unique biological labels. A fitted low-dimensional model can be useful for decoding even when its coordinates are not uniquely identifiable. Reviews of movement-related manifolds emphasize dynamical and task constraints \cite{Gallego2017}.

\section{Nonlinear latent models}
Sequential autoencoders and recurrent neural networks can infer nonlinear latent dynamics from single trials, under training objectives and architecture assumptions. Models such as LFADS use latent generative dynamics for population spike data. Flexible models can also overfit. Appropriate controls include held-out likelihood, behavior decoding, simulation-based identifiability tests, and stability across recording sessions.

\begin{workedexample}[A two-dimensional rotation]
Let $A=\begin{bmatrix}0&-\omega\\\omega&0\end{bmatrix}$. Then $\dot z=Az$ has solutions $z_1(t)=R\cos(\omega t+\phi)$ and $z_2(t)=R\sin(\omega t+\phi)$. The trajectory is a circle of radius $R$. A change of orthonormal coordinates rotates the picture but leaves its essential frequency unchanged.
\end{workedexample}
\begin{exercises}
\item Determine eigenvalues of the rotation matrix above.
\item Describe a dynamical system with a stable fixed point.
\item Explain why PCA dimensions and dynamical dimensions need not coincide.
\item Name two empirical checks required before claiming a neural manifold controls behavior.
\end{exercises}
\part{Scientific inference and reproducible methods}
\chapter{Hypothesis testing, resampling, and uncertainty}\label{ch:inference}
\begin{readergoal}Understand what a test proves, how uncertainty is quantified, and why a null model must preserve the right structure.\end{readergoal}
\section{What is a hypothesis test?}
A test begins by specifying a null family $H_0$ of possible data-generating mechanisms and a test statistic $T(D)$. The \emph{null distribution} describes the variation of $T$ under that family. A $p$-value is the probability, under the stipulated null model, of observing a statistic at least as extreme as the realized one according to the prespecified rule. It is not the posterior probability that the null is true, nor is it the probability the result occurred ``by chance'' without further qualification.

\section{Type I error, power, and effect size}
A level-$\alpha$ test controls the chance of rejecting a true null, provided its assumptions hold. Power is the rejection probability under a specified alternative. Statistical significance alone says little about effect magnitude. A coherence difference that is tiny but precisely estimated may be statistically significant, while a large effect in a small dataset may remain uncertain. Report the statistic, units, effect size, sampling unit, interval estimate, and model assumptions.

\section{Resampling methods}
A bootstrap approximates an estimator's sampling variability by resampling observational units from the data under specified conditions. The units might be animals or independent trials; sampling every LFP time point independently typically destroys serial dependence. A permutation test requires an exchangeability justification under the null, often by randomizing treatment labels under a randomized design. Surrogate time shifts can test certain synchrony claims only if their shift mechanism preserves the appropriate slow structure while breaking the tested fine-time association. These methods are not interchangeable.

\section{Hierarchical data and pseudoreplication}
Samples nested within a trial, trials within sessions, and sessions within animals are rarely independent. Treating every overlapping 100-ms analysis window as a fresh biological replicate inflates nominal sample size. A hierarchical model, cluster bootstrap, or analysis at the appropriate experimental unit can address some dependencies. Repeated statistical tests across channel pairs, frequency bins, or time bins require a multiple-comparison strategy adapted to the intended claim.

\section{Bias, generalization, and validation}
Preprocessing fitted on all the data before cross-validation can leak information from test folds into training. Selection of frequency bands after inspecting the outcome similarly biases later inference unless validated separately. Hold out whole animals or sessions when the scientific goal is cross-animal generalization. Predictive success on held-out data supports a model's useful regularities, not necessarily the literal truth of its latent biological explanation.

\begin{workedexample}[An exact conditional question]
To test whether two cells exhibit synchrony below a 5-ms scale beyond fixed 20-ms window spike counts, define windows, fixed spike counts, a synchrony statistic, and a resampling mechanism that randomizes fine timing conditional on the fixed counts. Any rejection should refer to the particular 20-ms coarse information held fixed, as emphasized in \cite{Amarasingham2012}.
\end{workedexample}
\begin{exercises}
\item State the difference between an effect-size estimate and a hypothesis-test decision.
\item Propose a hierarchical resampling unit for a dataset containing animals, sessions, trials, and time windows.
\item Explain why using the test set to choose a frequency band biases held-out performance.
\item Describe a null distribution that would be inappropriate for testing phase-amplitude coupling under strong signal nonstationarity.
\end{exercises}
\chapter{Numerical analysis and reproducible scientific computation}\label{ch:numerical}
\begin{readergoal}Translate mathematical equations into algorithms without overlooking finite precision, memory layout, and validation.\end{readergoal}
\section{A formula is not yet an algorithm}
A mathematical expression can define an exact quantity while computer code approximates it. Floating-point arithmetic rounds values and is not always associative: $(a+b)+c$ and $a+(b+c)$ can differ. An apparently exact bin count computed by $\lfloor T/\Delta\rfloor$ can be off by one if $T/\Delta$ is represented just below an integer. Implementations must state endpoint conventions and test boundary cases rather than rely on a universal ad hoc epsilon.

\section{Conditioning and algorithmic stability}
A well-posed problem has a solution that exists, is unique, and depends continuously on the data in an appropriate sense. An ill-conditioned problem may magnify small perturbations even under perfect implementation. A numerically unstable algorithm may magnify errors beyond what conditioning requires. Computing least-squares coefficients via QR or SVD is typically more stable than directly inverting $X^\top X$. Regularization stabilizes some inverse problems by choosing a constrained or penalized solution, at the cost of explicit modeling assumptions.

\section{Memory and complexity}
A dense sample covariance among $N$ neurons requires $O(N^2)$ storage, while a full matrix decomposition can be computationally expensive for large $N$. When sample count $T$ is much smaller than $N$, methods operating on a $T\times T$ Gram matrix can be more efficient. Signal operations can frequently be evaluated chunk by chunk, although filtering chunks requires state continuity or overlap handling. Memory mapping avoids loading full arrays but does not eliminate I/O costs.

\section{Testing the mathematical contract}
Good numerical tests use known synthetic inputs: a sinusoid with known frequency for PSD estimation; a complex exponential for phase operations; a Poisson simulation for count statistics; a low-rank signal with controlled noise for PCA; and a known VAR for Granger analysis. Property-based tests can verify constraints such as covariance symmetry, nonnegative variance, $\mathrm{PLV}\in[0,1]$, and consistency of concatenated block processing with full-array processing within a stated tolerance.

\section{Scientific metadata as part of the result}
An estimate needs sampling rate, time origin, measurement units, reference configuration, filtering parameters, electrode and trial identifiers, quality masks, random seeds, and software versions. Metadata are part of the mathematical map between the array and its real-world interpretation. A floating-point matrix detached from its coordinates can be mathematically valid yet scientifically unusable.

\begin{exercises}
\item Explain why zero-padding an FFT does not improve physical spectral resolution.
\item Design a test for chunked FIR filtering that detects boundary discontinuities.
\item List the metadata necessary to interpret a Hilbert-envelope moment series from 10-second windows.
\end{exercises}
\chapter{A complete path from first definitions to real datasets}\label{ch:roadmap}
\begin{readergoal}Practice the complete mathematical chain from a statement in English to a reproducible conclusion, with simulations preceding claims about physiology.\end{readergoal}
\section{Project I: two signals that look coupled but are not}
Start from the English claim: ``two channels have high coherence because the associated neurons communicate.'' Translate the measurable component to a statistic $\mathcal C_{xy}(f)$ and identify the unsupported causal clause. Construct independent Gaussian noises $n_x,n_y$ and a common oscillatory source $s$, then define $x=s+n_x$ and $y=s+n_y$. Calculate coherence, add measurement noise, and vary reference contamination. Compare estimates with and without the common source. Show that coherence may be high without direct communication. This project is motivated by \cite{Bastos2016,CohenTsuchiya2018}.

\section{Project II: what a theta peak actually means}
Simulate a process with a $1/f$-like background and a genuine narrowband theta oscillator. Next simulate only the aperiodic background, vary its exponent, and compute band power in both cases. Apply periodogram, Welch, and multitaper estimators; then compare parametric aperiodic fitting and an IRASA-type separation. Quantify false positives and biases across simulated repetitions before trying the methods on LFP. Here the hypotheses must distinguish a local peak \emph{over background} from integrated power over a fixed canonical band \cite{Donoghue2020,Gerster2022}.

\section{Project III: a manufactured phase-amplitude relationship}
Simulate a nonsinusoidal slow oscillator with harmonic content, and separately simulate a high-frequency oscillator whose amplitude truly depends on slow phase. Filter both into low and high bands, compute analytic signals, and compare PAC measures. Repeat with sharp transient waveforms and phase surrogates that preserve the important nuisance structure. Identify which simulated condition actually contains two distinct interacting generators. The validation target is the generation mechanism, not merely a significant PAC statistic \cite{Aru2015,ColeVoytek2018}.

\section{Project IV: spikes, jitter, and conditional inference}
Simulate independent Poisson spike trains sharing a slow time-dependent firing-rate modulation. A naive cross-correlogram may show coincidence peaks from rate cofluctuations. Define a conditional interval-jitter hypothesis holding coarse counts fixed and re-evaluate fine-time synchrony. Add precisely coordinated spikes and measure power across timescales and window widths. The lesson is to connect every resampling decision to an explicit null hypothesis \cite{Amarasingham2012}.

\section{Project V: latent neuronal dynamics}
Simulate $z_{t+1}=Az_t+w_t$ for a two-dimensional rotating latent state, then observe $N$ neurons through $r_t=Cz_t+\varepsilon_t$. Compare PCA, factor analysis, and a fitted linear state-space model. Vary the number of observed neurons, number of trials, observation noise, and dimensionality. Track held-out prediction and recovery of the latent subspace up to rotational ambiguity; avoid claiming exact recovery of nonidentifiable coordinate axes \cite{Churchland2012,CunninghamYu2014,Gallego2017}.

\section{A research conclusion requires four layers}
For each study, the final report should distinguish (1) the original English claim; (2) the formal mathematical target and its assumptions; (3) the numerical algorithm and validation outcomes; and (4) the physiological interpretation and its remaining alternatives. A clean implementation establishes only the third layer. A statistically significant result speaks to the second within its null model. A mechanistic conclusion requires evidence beyond either one in isolation.

\begin{exercises}
\item Draft a preregistered null hypothesis for one project, including the sampling unit, statistic, and exclusion rule.
\item Create a matrix matching five methods to their scientific estimands, key assumptions, and primary failure modes.
\item Design one negative control and one positive control for your chosen method.
\end{exercises}
\appendix
\chapter{Notation, assumptions, and units}
\begin{center}\begin{tabular}{p{0.29\textwidth}p{0.60\textwidth}}\toprule
Symbol & Meaning\\\midrule
$\in,\subseteq$ & Membership in a set; inclusion of one set in another\\
$\forall,\exists,\neg,\Rightarrow$ & For all; there exists; not; logical implication\\
$\mathbb N_0,\mathbb Z,\mathbb Q,\mathbb R,\mathbb C$ & Natural including zero, integer, rational, real, complex numbers\\
$\|x\|_2,\langle x,y\rangle$ & Euclidean norm and inner product\\
$E[X],\Var(X),\Cov(X,Y)$ & Expectation, variance, covariance\\
$\hat\theta,\theta$ & Estimator/estimate and model parameter\\
$f_s,\Delta t,N$ & Sampling frequency (Hz), interval (s), sample count\\
$x[n],x(t)$ & Discrete sampled signal and continuous-time signal\\
$X[k],S_{xx}(f)$ & DFT coefficient; power spectral density\\
$z(t),a(t),\phi(t)$ & Analytic signal, amplitude envelope, phase\\
$N(t),\lambda(t\mid\mathcal H_t)$ & Point-process count and conditional intensity\\
$\mathcal C_{xy}(f)$ & Magnitude-squared coherence\\
$X\in\mathbb R^{T\times N}$ & Population matrix: times by neurons\\
$\mathcal X\in\mathbb R^{N\times T\times R}$ & Tensor: neurons by time by trials\\
\bottomrule\end{tabular}\end{center}

\section*{Explicit conventions}
PSD values have signal-units-squared per Hz under a density convention; firing rate uses spikes/s; phase is measured in radians unless labeled otherwise. Distinguish the sample count $N$ from the point-process count $N(t)$ using context. The symbol $T$ may denote a recording duration in signal chapters but number of observations in matrix expressions when explicitly defined. We use $\log$ for natural logarithm. Complex inner products are linear in their first argument by the convention stated in Chapter~\ref{ch:functionspaces}. A reported result must disclose its observational unit and any filtering or binning convention.
\chapter{Provisional mathematical map to \texttt{arsenal/path}}\label{app:arsenal}
This is a \textbf{conceptual map, not a verified file inventory}. The GitHub repository and the actual \texttt{arsenal/path} tree were not available during compilation. Names below describe analysis families and do not assert that particular MATLAB files are implemented.
\begin{longtable}{p{0.25\textwidth}p{0.34\textwidth}p{0.31\textwidth}}\toprule
Method family & Mathematical prerequisites & Checks before interpreting results\\\midrule\endhead
Binning and firing rates & Natural numbers; functions; point processes & Window edges, exposure, refractory period\\
FFT and PSD & Complex numbers; sequences; convolution & Sample rate, leakage, normalization, 1/f background\\
Hilbert envelopes & Complex phase; Fourier; filtering & Narrowband validity, edges, amplitude near zero\\
Sliding moments & Probability; expectation; central moments & Window overlap, variance conventions, autocorrelation\\
Coherence and PLV & Cross-spectra; circular statistics & Reference, common input, sample size, null design\\
PAC & Analytic signals; nonlinear statistics & Harmonics, transients, nonstationarity, surrogates\\
Granger and VAR & Linear algebra; stochastic processes; likelihood & Lag order, stationarity, latent common input\\
Spike synchrony & Point processes; conditional probability & Coarse counts, jitter window, precise null\\
PCA and SVD & Linear maps; covariance; spectral theorem & Scaling, noise, held-out stability\\
State-space models & EDOs; conditional inference; linear algebra & Observability, latent-coordinate ambiguity\\
CSD and inverse models & Multivariable derivatives; forward physics & Geometry, conductivity, regularization\\
\bottomrule\end{longtable}
For a future verified edition, each row should link to the exact MATLAB function, an explicit function signature, the corresponding mathematical equation, a reference, a tested example, and at least one synthetic negative control. An automatically generated index from the real repository would be preferable to copying guessed filenames.
\chapter{Selected answers, proof sketches, and further problems}
\section{Selected answers}
\paragraph{Why empty collections matter (Chapter~\ref{ch:english}).} The statement $\forall x\in\varnothing\,P(x)$ is true because no element provides a counterexample. This is not evidence that any spike occurred: it is a logical consequence of quantifying over an empty domain.
\paragraph{Conditional order (Chapter~\ref{ch:logic}).} ``Every cell fires at some time'' has the form $\forall u\exists tF(u,t)$. ``There exists a time when every cell fires'' is $\exists t\forall uF(u,t)$. The second implies the first for a nonempty common time domain, but not conversely.
\paragraph{Count normalization (Chapter~\ref{ch:rationals}).} Twenty spikes in 5 s gives 4/s; thirty spikes in 15 s gives 2/s. Their unweighted mean is 3/s; the pooled rate is $50/20=2.5$/s. The different answers correspond to different weighting rules.
\paragraph{Projection (Chapter~\ref{ch:linear}).} $P^2=UU^\top UU^\top=U I U^\top=P$. Also $P^\top=P$. For any $v$ in the column span of $U$, $Pv=v$.
\paragraph{Gain model (Chapter~\ref{ch:population}).} Because $N\mid G$ is Poisson with mean $G\mu$, total variance yields $\Var N=E[G\mu]+\Var(G\mu)=\mu+\mu^2\Var G$ when $E G=1$.
\section{Capstone proof and analysis problems}
\begin{enumerate}
\item Start from the definition of a function and develop the meaning of a sampled time series, including its domain, range, units, and missing data.
\item Prove that the covariance matrix of a finite centered real data matrix is positive semidefinite. Identify when it fails to be positive definite.
\item From a linear state-space model, construct two different latent coordinate systems that generate identical observations.
\item Compare a phase-locking measure with an amplitude-weighted PAC measure on signals with very low amplitude. Characterize instability.
\item Write a full conditional resampling test for spike synchrony, explicitly stating what remains fixed under the null.
\item Derive the memory footprint and main arithmetic complexity of a covariance-based versus an SVD-based approach for $T\ll N$.
\item Design and report a blinded synthetic benchmark where an algorithm is required to distinguish an oscillation from a steep aperiodic spectrum.
\end{enumerate}
These problems are deliberately cumulative. They require assembling definitions from several parts rather than treating methods as isolated formulas.
\backmatter
\begin{thebibliography}{99}
\bibitem{Amarasingham2012} Amarasingham, Asohan and Harrison, Matthew T. and Hatsopoulos, Nicholas G. and Geman, Stuart. \emph{Conditional modeling and the jitter method of spike resampling}. Journal of Neurophysiology, 2012. \url{https://doi.org/10.1152/jn.00633.2011}
\bibitem{Aru2015} Aru, Juhan and Aru, Jaan and Priesemann, Viola and others. \emph{Untangling cross-frequency coupling in neuroscience}. Current Opinion in Neurobiology, 2015. \url{https://doi.org/10.1016/j.conb.2014.08.002}
\bibitem{Bastos2016} Bastos, Andre M. and Schoffelen, Jan-Mathijs. \emph{A Tutorial Review of Functional Connectivity Analysis Methods and Their Interpretational Pitfalls}. Frontiers in Systems Neuroscience, 2016. \url{https://doi.org/10.3389/fnsys.2015.00175}
\bibitem{Berens2009} Berens, Philipp. \emph{CircStat: A MATLAB Toolbox for Circular Statistics}. Journal of Statistical Software, 2009. \url{https://doi.org/10.18637/jss.v031.i10}
\bibitem{Bokil2010} Bokil, Hemant and Andrews, Peter and Kulkarni, Jayant E. and Mehta, Samar and Mitra, Partha. \emph{Chronux: A Platform for Analyzing Neural Signals}. Journal of Neuroscience Methods, 2010. \url{https://doi.org/10.1016/j.jneumeth.2010.06.020}
\bibitem{Buccino2020} Buccino, Alessio P. and others. \emph{SpikeInterface, a unified framework for spike sorting}. eLife, 2020. \url{https://doi.org/10.7554/eLife.61834}
\bibitem{Buccino2026} Buccino, Alessio Paolo and Sridhar, Arjun and Feng, David and Svoboda, Karel and Siegle, Joshua H.. \emph{Efficient and reproducible pipelines for spike sorting large-scale electrophysiology data}. eLife, 2026. \url{https://doi.org/10.7554/eLife.110170}
\bibitem{Buzsaki2012} Buzsaki, Gyorgy and Anastassiou, Costas A. and Koch, Christof. \emph{The origin of extracellular fields and currents---EEG, ECoG, LFP and spikes}. Nature Reviews Neuroscience, 2012. \url{https://doi.org/10.1038/nrn3241}
\bibitem{Cekic2018} Cekic, Sezen and Grandjean, Didier and Renaud, Olivier. \emph{Time, frequency, and time-varying Granger-causality measures in neuroscience}. Statistics in Medicine, 2018. \url{https://doi.org/10.1002/sim.7621}
\bibitem{Churchland2012} Churchland, Mark M. and others. \emph{Neural population dynamics during reaching}. Nature, 2012. \url{https://doi.org/10.1038/nature11129}
\bibitem{CohenKohn2011} Cohen, Marlene R. and Kohn, Adam. \emph{Measuring and interpreting neuronal correlations}. Nature Neuroscience, 2011. \url{https://doi.org/10.1038/nn.2842}
\bibitem{CohenTsuchiya2018} Cohen, Dror and Tsuchiya, Naotsugu. \emph{The Effect of Common Signals on Power, Coherence and Granger Causality}. Frontiers in Systems Neuroscience, 2018. \url{https://doi.org/10.3389/fnsys.2018.00030}
\bibitem{ColeVoytek2018} Cole, Scott and Voytek, Bradley. \emph{Cycle-by-cycle analysis of neural oscillations}. bioRxiv, 2018. \url{https://doi.org/10.1101/302000}
\bibitem{CheveigneNelken2019} de Cheveigne, Alain and Nelken, Israel. \emph{Filters: When, Why, and How (Not) to Use Them}. Neuron, 2019. \url{https://doi.org/10.1016/j.neuron.2019.02.039}
\bibitem{Donoghue2020} Donoghue, Thomas and others. \emph{Parameterizing neural power spectra into periodic and aperiodic components}. Nature Neuroscience, 2020. \url{https://doi.org/10.1038/s41593-020-00744-x}
\bibitem{Einevoll2013} Einevoll, Gaute T. and Kayser, Christoph and Logothetis, Nikos K. and Panzeri, Stefano. \emph{Modelling and analysis of local field potentials for studying the function of cortical circuits}. Nature Reviews Neuroscience, 2013. \url{https://doi.org/10.1038/nrn3599}
\bibitem{Gallego2017} Gallego, Juan A. and Perich, Matthew G. and Miller, Lee E. and Solla, Sara A.. \emph{Neural Manifolds for the Control of Movement}. Neuron, 2017. \url{https://doi.org/10.1016/j.neuron.2017.05.025}
\bibitem{Gao2017} Gao, Peiran and others. \emph{A theory of multineuronal dimensionality, dynamics and measurement}. bioRxiv, 2017. \url{https://doi.org/10.1101/214262}
\bibitem{Gerster2022} Gerster, Moritz and others. \emph{Separating Neural Oscillations from Aperiodic 1/f Activity: Challenges and Recommendations}. Neuroinformatics, 2022. \url{https://doi.org/10.1007/s12021-022-09581-8}
\bibitem{Goris2014} Goris, Robbe L. T. and Movshon, J. Anthony and Simoncelli, Eero P.. \emph{Partitioning neuronal variability}. Nature Neuroscience, 2014. \url{https://doi.org/10.1038/nn.3711}
\bibitem{CunninghamYu2014} Cunningham, John P. and Yu, Byron M.. \emph{Dimensionality reduction for large-scale neural recordings}. Nature Neuroscience, 2014. \url{https://doi.org/10.1038/nn.3776}
\bibitem{HalmosNaive} Halmos, Paul R.. \emph{Naive Set Theory}. Springer, 1974. \url{https://doi.org/10.1007/978-1-4757-1645-0}
\bibitem{MendelsonNumbers} Mendelson, Elliott. \emph{Number Systems and the Foundations of Analysis}. Dover Publications, 2008. 
\bibitem{Canolty2006} Canolty, R. T. and others. \emph{High Gamma Power Is Phase-Locked to Theta Oscillations in Human Neocortex}. Science, 2006. \url{https://doi.org/10.1126/science.1128115}
\end{thebibliography}
\end{document}